% =====================================================================================
% TFM — Predicción de la degradación por ruido en circuitos cuánticos NISQ antes de su
%       ejecución: comparativa entre modelos tabulares y redes neuronales de grafos
% Gregory Harutyunyan Gevorgyan · Máster en Inteligencia Artificial · UNIR
%
% Estructura y estado de cada capítulo: doc/memoria/estructura.md
% Hallazgos y cifras, con su origen:     doc/memoria/materiales.md
% Qué falta y qué NO se puede citar aún: doc/memoria/pendientes.md
% =====================================================================================
\documentclass[11pt,a4paper,spanish]{book}

% --- Estilo institucional de la UNIR ----------------------------------------------
% ⚠️ CALCADO DE LA PLANTILLA OFICIAL (sep-2026), por decisión del autor. La plantilla
%    carga SOLO `estilo_unir-1`, que a su vez trae babel e inputenc sin opciones; el
%    idioma llega por la opción `spanish` de \documentclass. Se retiraron por tanto los
%    \usepackage de inputenc, fontenc y babel que la memoria cargaba antes del estilo.
\usepackage{estilo_unir-1}

% --- El resto del preámbulo -------------------------------------------------------
% --- Codificación y Lenguaje ---
\usepackage[utf8]{inputenc}
\usepackage[T1]{fontenc}
\usepackage[spanish,es-tabla]{babel}

% --- Geometría y Diseño ---
\usepackage[margin=2.5cm]{geometry}
\usepackage{setspace}
\onehalfspacing % Interlineado 1.5

% --- Matemáticas Avanzadas (Crucial para Cuántica) ---
\usepackage{amsmath, amsfonts, amssymb, amsthm}
\usepackage{physics} % Proporciona macros para brakets: \bra{}, \ket{}, \braket{}
\usepackage{bm}      % Negritas en símbolos matemáticos

% --- Gráficos y Tablas ---
\usepackage{graphicx}
\usepackage{booktabs}
\usepackage{xcolor}
\usepackage{hyperref}
\hypersetup{
    colorlinks=true,
    linkcolor=blue,
    filecolor=magenta,      
    urlcolor=cyan,
    citecolor=black
}

% --- Algoritmos y Código ---
\usepackage{listings}
\definecolor{codegreen}{rgb}{0,0.6,0}
\lstset{
    commentstyle=\color{codegreen},
    keywordstyle=\color{magenta},
    numberstyle=\tiny\color{gray},
    stringstyle=\color{purple},
    basicstyle=\ttfamily\small,
    breaklines=true,
    captionpos=b,
    keepspaces=true,
    showspaces=false,
    language=Python
}

\usepackage{apacite}


\let\rank\undefined % Forzamos que \rank quede libre

\numberwithin{equation}{chapter}
%\usepackage{caption}
\numberwithin{figure}{chapter}
%\usepackage{chngcntr}
%\counterwithin{equation}{chapter}
%\counterwithin{figure}{chapter}
%\counterwithin{table}{chapter}
%\renewcommand\theequation{\thechapter.\arabic{equation}}
%\renewcommand\thefigure{\thechapter.\arabic{figure}}
%\renewcommand\thetable{\thechapter.\arabic{table}}
%\renewcommand\theequation{\thechapter.\arabic{equation}}
%\counterwithin{figure}{chapter}
%\counterwithin{table}{chapter}
%\renewcommand\thefigure{\thechapter.\arabic{figure}}
%\renewcommand\thetable{\thechapter.\arabic{table}}
%\renewcommand\thefigure{\thechapter.\arabic{figure}}
%\makeatletter
%\renewcommand\p@figure{\thechapter..\arabic{figure}}
%\makeatother


% =====================================================================================
% Portada
% =====================================================================================
% ⚠️ TITULO CAMBIADO (sep-2026) por decision del autor. El anterior era «Mitigación de
%    Ruido Cuántico en Hardware NISQ mediante Deep Learning»; se retiró porque anunciaba
%    como método el que NO gana la comparativa y omitía el diferenciador del trabajo, que
%    es predecir ANTES de ejecutar. PENDIENTE: confirmarlo con el director.
% ⚠️ SALTOS DE LINEA EXPLICITOS con \newline (NO \\: dentro de la tabular del estilo, \\
%    cerraria la fila y descolocaria la portada entera).
%    Se fijan a mano porque la columna p{12cm} del estilo justifica el titulo a cuerpo \Huge,
%    y a ~26 caracteres por linea eso deja huecos enormes entre palabras. Cada \newline
%    cierra su linea sin estirarla. Los cortes son semanticos, no arbitrarios.
%    ⚠️ SI SE CAMBIA EL TITULO, HAY QUE REHACER LOS CORTES.
\title{Predicción de la\newline
degradación por ruido en\newline
circuitos cuánticos NISQ\newline
antes de su ejecución:\newline
comparativa entre modelos\newline
tabulares y redes\newline
neuronales de grafos}
\titulacion{Máster Universitario en Inteligencia Artificial}
\author{Gregory Harutyunyan Gevorgyan}
\director{Pablo Abellán Galiana}
\nombreciudad{ Granada}
% ⚠️ Fecha de entrega, la misma que fija el criterio «T» del objetivo general (§4.1).
%    NO usar \today: la portada quedaría con la fecha de la ultima compilacion.
\date{16 de agosto de 2026}

% =====================================================================================
\begin{document}

\renewcommand{\listfigurename}{Índice de Ilustraciones}
\renewcommand{\listtablename}{Índice de Tablas}
\renewcommand{\contentsname}{Índice de Contenidos}
\renewcommand{\figurename}{Figura}
\renewcommand{\tablename}{Tabla}

% ⚠️ EL TITULO NO SE PARTE CON GUION EN LA PORTADA.
%    Los parametros van AQUI, envolviendo a \maketitle, y NO dentro de \title{}: el estilo
%    compone el titulo dentro de grupos anidados que se cierran antes de romper el parrafo,
%    asi que puestos ahi dentro se revierten y no surten efecto. Sin esto la portada parte
%    «degrada-ción», «en-tre» y «re-des», que a cuerpo \Huge canta mucho.
%    \raggedright evita ademas que la justificacion abra huecos enormes entre palabras.
%    El estilo institucional (estilo_unir-1.sty) NO se toca: es calco de la plantilla.
\begingroup
\hyphenpenalty=10000 \exhyphenpenalty=10000 \raggedright
\maketitle
\endgroup

% ⚠️ NO se usa \frontmatter tal cual: su \cleardoublepage empujaba a pagina impar y dejaba
%    una HOJA EN BLANCO detras de la portada (la antigua pagina 2). Aqui se hace lo mismo
%    que hace \frontmatter —numeracion romana y capitulos sin numerar— pero sin ese salto.
%    \chapter ya no fuerza pagina impar por su cuenta: estilo_unir-1.sty lo parchea (l.50).
\makeatletter
\@mainmatterfalse
\makeatother
\pagenumbering{roman}

% -------------------------------------------------------------------------------------
% ⚠️ ORDEN FIJADO POR EL AUTOR (sep-2026): Resumen y Abstract van JUSTO DETRAS de la
%    portada y ANTES de los indices. No devolverlos detras del indice.
% -------------------------------------------------------------------------------------
\chapter{Resumen}
% ⚠️ EXTENSIÓN: las Instrucciones §1.2 dicen «máximo 150 palabras» y las §2.1 «de 150 a 300».
%    Los dos documentos oficiales se contradicen. EL AUTOR DECIDIO (sep-2026) acogerse a la
%    horquilla de §2.1 para poder incluir la justificacion y las contribuciones. MEDIDO: 241.
%    Si se reescribe, volver a contarlas y NO pasar de 300.
%
% ⚠️ PALABRAS CLAVE: NO poner «graph transformer». Los modelos de este trabajo son MPNN
%    (SAGEConv, GINConv, GATv2Conv); un transformer de grafo es otra arquitectura y aquí
%    no se usa ninguna. Se corrigió en sep-2026.

Los ordenadores cuánticos actuales devuelven resultados degradados por el ruido. Las técnicas
habituales para corregirlo actúan después de medir y exigen ejecutar el circuito muchas veces,
lo que multiplica el consumo de un recurso caro y escaso: el tiempo de procesador cuántico.
Anticipar esa degradación permitiría decidir qué versión de un circuito merece la pena
ejecutar, o si merece la pena ejecutarlo siquiera, sin gastar ese tiempo. Este trabajo estudia
si la degradación puede predecirse antes de ejecutar, solo a partir de la estructura del
circuito y de la calibración del procesador.

Para ello se construye un conjunto de veinte mil circuitos de cinco a quince qubits,
etiquetados con su degradación exacta sobre la calibración real de un procesador de IBM
durante cuarenta y dos días, y particionado en tres ejes de generalización independientes:
tipo de circuito, tamaño y tiempo. Sobre él se comparan tres tipos de modelos bajo un mismo protocolo: 
una regresión Ridge, un Random Forest y una red neuronal de grafos.

El trabajo aporta tres resultados. Primero, la degradación con signo no es predecible antes de
ejecutar y su reformulación como factor
de supervivencia sí lo es. Segundo, Random Forest gana la comparativa, y la representación en
grafo solo compensa al extrapolar a circuitos mayores o a algoritmos nunca vistos. Tercero,
acertar sobre el factor no garantiza mitigar bien: Ridge lo predice con coeficiente positivo
en los cinco observables y aun así empeora cuatro de ellos al aplicarlo.

\textbf{Palabras clave:} computación cuántica, NISQ, mitigación de errores, aprendizaje
automático, redes neuronales de grafos.

% -------------------------------------------------------------------------------------
\chapter{Abstract}
% Traducción del resumen en español. Si se toca uno, tocar el otro. MEDIDO: 231 palabras.

Today's quantum computers return results degraded by noise. The usual techniques for
correcting it act after measurement and require running the circuit many times, which
multiplies the consumption of an expensive and scarce resource: quantum processor time.
Anticipating that degradation would make it possible to decide which version of a circuit is
worth running, or whether it is worth running at all, without spending that time. This work
studies whether the degradation can be predicted before execution, using only the structure of
the circuit and the calibration of the processor.

To that end, a dataset of twenty thousand circuits of five to fifteen qubits is built,
labelled with their exact degradation over the real calibration of an IBM processor across
forty-two days, and partitioned along three independent axes of generalisation: circuit type,
size and time. Three types of model are then compared under a single protocol: a Ridge
regression, a Random Forest and a graph neural network.

The work yields three results. First, the signed degradation is not predictable before
execution, whereas its reformulation as a survival factor is. Second, Random Forest wins the comparison, and the graph representation
only pays off when extrapolating to larger circuits or to previously unseen algorithms. Third,
predicting the factor well does not guarantee good mitigation: Ridge predicts it with a
positive coefficient across all five observables and yet worsens four of them when applied.

\textbf{Keywords:} quantum computing, NISQ, error mitigation, machine learning, graph neural
networks.

% -------------------------------------------------------------------------------------
\tableofcontents
\listoffigures
\listoftables

% =====================================================================================
\mainmatter

% =====================================================================================
% Estructura: tipo de trabajo 3 — COMPARATIVA DE SOLUCIONES.
% Cada bloque corresponde a un apartado de la estructura oficial (Instrucciones §2.6).
% =====================================================================================

% --- Introducción -------------------------------------------------------------------
\chapter{Introducción}
\label{cap:intro}

% GUION. Fuentes: doc/memoria/estructura.md, doc/00_Justificaciones.md, doc/03_Ruido_Cuantico.md
% EXTENSION: 3-5 paginas (Instrucciones seccion 2.3). No mas.
%
% ⚠️ ESTRUCTURA FIJADA POR LA NORMA (seccion 2.3): Motivacion · Planteamiento del
%    problema · Estructura del trabajo. Aqui los objetivos van «a grandes rasgos»; los
%    objetivos formales, SMART, viven en el capitulo 5.
%
% ⚠️ DECISION DE REDACCION (sep-2026): la introduccion ANTICIPA el hallazgo negativo.
%    Motivo: §1.3 ya lo enuncia, y plantear §1.2.2 como pregunta abierta dejaba el
%    capitulo contradiciendose consigo mismo veinte lineas despues. Ademas, en un
%    trabajo de tipo 3 el tribunal espera saber desde el principio que se compara y
%    que salio.
%
% ⚠️ REVISAR AL CERRAR RESULTADOS: las contribuciones solo se enuncian del todo bien
%    cuando se sabe que se ha demostrado.

% Parrafo de entrada del capitulo. Solo dice QUE hay en cada apartado; el contenido no se
% adelanta aqui, que para eso estan los apartados.
En este capítulo se introduce el próposito de esta memoria. El apartado~\ref{sec:motivacion} explica la utilidad de un ordenador cuántico y cómo el ruido les perjudica. El
apartado~\ref{sec:planteamiento} plantea el problema concreto que este trabajo aborda: corregir el ruido
exige ejecutar el circuito, y ejecutarlo es caro; y enuncia la propuesta. Finalmente, el
apartado~\ref{sec:estructura} describe la organización de la memoria.

\section{Motivación}
\label{sec:motivacion}
% ⚠️ FUENTES RESTRINGIDAS POR DECISION DEL AUTOR (sep-2026): esta seccion 1.1 se apoya
%    UNICAMENTE en tres fuentes divulgativas:
%      1.1.1 -> ibm2026computacion   (IBM Think, «¿Que es la computacion cuantica?»)
%      1.1.2 -> tqi2026nisq          (The Quantum Insider, guia NISQ)
%      1.1.3 -> duque2025ruido       (Manu Duque, «Ruido cuantico»)
%    NO añadir aqui Shor, Grover, Preskill, Kim, Cai ni Nielsen: se citan mas adelante.


La computación cuántica es un campo de la informática y la ingeniería que aprovecha
las cualidades propias de la mecánica cuántica para resolver problemas que quedan fuera del
alcance incluso de los ordenadores clásicos más potentes \cite{ibm2026computacion}. Y esta definición
hay que leerla con cuidado porque para la mayoría de las tareas se espera que los ordenadores clásicos sigan siendo la mejor
solución. La ventaja cuántica no es general, sino que está
concentrada en un tipo muy concreto de problema.

Ese tipo de problema es aquel en el que muchas variables interactúan de forma complicada. 
Un ejemplo es modelar el
comportamiento de los átomos de una molécula, complejo por las múltiples interacciones entre
sus electrones \cite{ibm2026computacion}. Un ordenador clásico necesita
\emph{calcular} el comportamiento de la molécula, mientras que un ordenador cuántico no
necesita hacer los cálculos, porque puede \textbf{imitar} el sistema molecular.

% ⚠️ FORMATO: titulo arriba y fuente abajo (Instrucciones §1.3, p. 6).
% FOTOGRAFIA DE TERCEROS. Origen: la URL que aporto el autor,
%   https://i0.wp.com/clipset.com/wp-content/uploads/2018/01/ibm-cuantico4.jpg
% ⚠️ SIN LICENCIA DECLARADA en la fuente. Casi con seguridad es material grafico de IBM
%    republicado por el medio, o una foto tomada en la feria. Antes de la entrega hay que
%    localizar el origen primario y su licencia, o sustituirla por una del banco de
%    imagenes de IBM Research, que si publica condiciones de uso.
% ⚠️ Es una JPG de 960x1280 px. A 0.35\textwidth (5,3 cm) da ~460 dpi, de sobra para los
%    300 dpi que exige el README de figuras. NO ampliarla mas de 0.6\textwidth.
\begin{figure}[htbp]
    \centering
    \includegraphics[width=0.35\textwidth]{cap01/ibm_cuantico}

    \caption{Un procesador cuántico de IBM instalado en su criostato.}
    \label{fig:ibm_cuantico}

    {\footnotesize Fuente: Clipset (2018).}
\end{figure}

Lo que se ve en la Figura~\ref{fig:ibm_cuantico} no es el ordenador cuántico, sino el
frigorífico que lo contiene. El procesador cuántico es una
pieza del tamaño de una moneda que va sujeta abajo del todo, a una temperatura cercana al
cero absoluto. Toda esa maquinaria existe para un solo motivo: aislar los qubits del entorno,
porque cualquier vibración o cualquier resto de calor los perturba. Esa fragilidad es
justamente el problema del que trata esta memoria.

En general para hacerse una idea, los ordenadores cuánticos sirven, o se prevé que servirán, para solventar problemas
complejos en campos como: industria farmacéutica, química y ciencia de materiales, energía, finanzas, biología computacional, aprendizaje automático, ciberseguridad y semiconductores \cite{ibm2026computacion}.  Históricamente, la primera aplicación
práctica que se propuso fue el algoritmo de factorización de enteros
publicado por Peter Shor en 1994 \cite{shor1994algorithms}, que mostró cómo una máquina cuántica podría llegar a
romper sistemas criptográficos.

Sin embargo, aunque la computación cuántica sea una tecnología tan prometedora, todavía presenta
ciertos desafíos y limitaciones importantes. 

Actualmente, todos los ordenadores cuánticos disponibles comercialmente (los de IBM, Google, IonQ\dots) pertenecen a una misma categoría, la llamada \textbf{era NISQ}, \textit{Noisy Intermediate-Scale Quantum} \cite{tqi2026nisq}. El físico John Preskill acuñó este término en 2018 para describir una generación de ordenadores cuánticos que ya son difíciles de simular con ordenadores clásicos, pero que siguen teniendo demasiados errores como para aprovechar todo el potencial en aplicaciones reales.
\cite{tqi2026nisq}.

Los qubits, la unidad básica de información en la computación cuántica, son extremadamente sensibles a cualquier perturbación de su entorno, y cada operación introduce una pequeña posibilidad de error. Estos errores se acumulan con el tiempo y limita la fiabilidad de los resultados obtenidos. Lo actuales sistemas disponen de entre decenas y unos pocos miles de qubits que no son suficientes para que lo ordenadores cuánticos sean totalmente fiables. \cite{tqi2026nisq}.

Aunque se conoce una posible solución, llamada Correción de Errores Cuánticos (QEC, por sus siglas en inglés \textit{Quantum Error Correction}), llevarla a la práctica sigue siendo uno de los mayores desafíos de la computación cuántica. Para hacerse una idea: para conseguir una buena fidelidad de un solo qubit es necesario ejecutar cientos o miles de ellos \cite{tqi2026nisq}. Esto hace que los ordenadores cuánticos actuales estén aún lejos de poder aplicarse de manera eficaz.

Mientras tanto, existe una alternativa práctica en la que se centra nuestro trabajo: la Mitigación de Errores Cuánticos (QEM, por sus siglas en inglés \textit{Quantum Error Correction}). En lugar de intentar eliminar el ruido, su objetivo es estimar cuál habría sido el resultado obtenido en un escenario ideal sin errores. Algunas de las técnicas que se incluyen en este área son: \textit{Zero Noise Extrapolation} (ZNE), que estima el resultado ideal extrapolando las mediciones obtenidas de ruido; \textit{Probabilistic Error Cancellation} (PEC), que compensa los efectos del ruido mediante una combinación probabilística; y \textit{Readout Error Mitigation} (REM), que corrige los errores producidos durante la etapa de medición de los qubits..


\section{Planteamiento del problema}
\label{sec:planteamiento}

\subsection{Limitaciones de los enfoques actuales QEM}
\label{sec:limitaciones}

El problema de estos métodos no es la precisión, sino el \textbf{coste}. Tanto ZNE como PEC
necesitan ejecutar el circuito \textbf{varias veces más} para poder estimar cuánto lo ha desviado el ruido. Como consecuencia, 
la corrección de errores necesita una 
cantidad adicional de tiempo de ejecución en el procesador cuántico, un recurso escaso y caro en la actualidad.

De ahí el interés, en este ámbito, del \textbf{aprendizaje automático} (ML-QEM, \textit{Machine Learning - Quantum Error Mitigation}): si un modelo aprende la relación entre el ruido y su efecto, la corrección se calcula sin gastar ejecuciones adicionales. El procedimiento general consta de tres pasos \cite{liao2024machine}:

\begin{enumerate}
      \item \textbf{Generación.} Se produce una gran cantidad de circuitos de entrenamiento
          que imitan la estructura del circuito de interés, de forma que el modelo aprenda
          sobre ejemplos parecidos a aquel que después tendrá que corregir.
      \item \textbf{Codificación.} Cada circuito se pasa por un codificador que extrae sus
          características (el número de compuertas, su tipo o su profundidad, entre otras),
          convirtiendo el circuito en una representación que el modelo pueda consumir.
      \item \textbf{Etiquetado.} Los circuitos se ejecutan en la QPU (Quantum Processing Unit) para obtener sus valores esperados con ruido. La comparación con el valor ideal conocido permite calcular la desviación introducida por el ruido, que se utiliza como variable objetivo del modelo.
\end{enumerate}

Con ese conjunto de datos el modelo aprende a estimar la corrección, y una vez entrenado la
aplica sin necesidad de las ejecuciones extra que hacen falta en ZNE o PEC. Ahora bien, el esquema
tiene una limitación: 

\begin{itemize}
      \item el modelo toma como entrada el \textbf{valor ruidoso ya medido}, y por tanto \textbf{es necesario haber ejecutado el circuito} para poder corregirlo. La corrección sigue llegando después del gasto,
      \item y el \textbf{valor ideal} debe conocerse de antemano para entrenar el modelo.
\end{itemize}

% ⚠️ FORMATO: titulo arriba y fuente abajo (Instrucciones §1.3, p. 6).
% Figura propia. Fuente editable en figuras/cap01/mlqem_flujo.svg; el PDF se regenera con
%   rsvg-convert -f pdf -o mlqem_flujo.pdf mlqem_flujo.svg
% Representa la enumeracion de esta misma seccion Y la limitacion que la sigue: el trazo
% rojo discontinuo es el valor ruidoso medido entrando como entrada del modelo, que es
% exactamente lo que obliga a ejecutar el circuito antes de poder corregirlo. Sin ese
% trazo la figura seria una ilustracion redundante del texto de al lado.
\begin{figure}[htbp]
    \centering
    \includegraphics[width=0.92\textwidth]{cap01/mlqem_flujo}

    \caption{El esquema general de la mitigación por aprendizaje automático.}
    \label{fig:mlqem}

    {\footnotesize Fuente: Claude Code \cite{anthropic_claude}.}
\end{figure}

Para conseguir el valor ideal existen varios métodos pero se ejecutan en ordenadores clásicos. Como mostraremos más adelante, los ordenadores clásicos son incapaces de reproducir los procesos de los ordenadores cuánticos a partir de un número de qubits.

\subsection{La pregunta de este trabajo}
\label{sec:pregunta}

De ahí surge la pregunta que se trata en la memoria:

\begin{center} 
      \doublebox{
      \begin{minipage}{0.85\linewidth} \vspace{0.5em}
             \centering 
             \textit{¿Se puede predecir la degradación que sufrirá un circuito \textbf{antes de ejecutarlo}, mirando únicamente su estructura y la calibración del procesador?} \vspace{0.5em} \end{minipage}} 
\end{center}

Lo que se propone no es cambiar el método, sino el momento en que se aplica la corrección (véase el Capítulo~\ref{sec:reformulacion}): en lugar de corregir el resultado una vez obtenido, intentar \textbf{anticiparlo}. Si esto fuera posible, su utilidad sería inmediata y diferente a la de los métodos anteriores. Permitirá estimar de antemano qué versión de un circuito merece la pena ejecutar, o incluso si merece la pena ejecutarlo, sin consumir tiempo de procesador cuántico. Aunque la literatura ha explorado ideas similares, como clasificar circuitos según la calidad esperada de sus resultados antes de ejecutarlos \cite{hartnett2024learning}, todavía el enfoque `evitando recursos cuánticos' sigue contando con escasas contribuciones. 

\subsection{Propuesta}
\label{sec:propuesta}

Responder a esa pregunta exige, antes que un modelo, un conjunto de datos que no existe.
De modo que lo primero que se propone es construirlo: \textbf{veinte mil circuitos} de entre
cinco y quince qubits, simulados sobre el gemelo digital de un procesador de IBM y
etiquetados con la degradación \textbf{exacta} que sufre cada uno. El ruido no se inventa, se
reconstruye día a día a partir de la \textbf{calibración real} del procesador durante cuarenta y
dos días. Y la partición se diseña para que el conjunto de prueba obligue a
generalizar en \textbf{tres ejes independientes} (el tipo de circuito, su tamaño y el día de
calibración).

Sobre esos datos se comparan \textbf{tres soluciones} bajo un mismo protocolo: una regresión
lineal regularizada (\textbf{Ridge}), un \textbf{Random Forest} (el mejor modelo tabular de referencia
según la ablación de \citeA{liao2024machine}) y una \textbf{red neuronal de grafos con paso de
mensajes} (MPNN, \textit{message passing neural network}). La elección no es arbitraria,
porque un circuito no es una fila: es una matriz de altura variable, una fila por puerta.
Los dos primeros modelos exigen filas de longitud fija, así que hay que \textbf{resumir} el
circuito en columnas, y ese resumen tira el orden y la estructura (dos circuitos
con las mismas puertas estructuradas de formas distintas generan la misma fila). Sin embargo, la red de grafos consume el circuito tal cual. La
diferencia entre unos y otra es, por tanto, \textbf{la medida de cuánto vale la estructura que el
resumen descarta}. Este planteamiento está respaldado por trabajos previos. Por ejemplo, \citeA{tudisco2025graph} demuestra, antes de introducir su red basada en grafos, que características globales como el número de qubits, la profundidad del circuito o el recuento de puertas no son suficientes para distinguir adecuadamente las clases consideradas. Además, la comparación entre redes de grafos y Random Forest aplicada a circuitos cuánticos ya cuenta con un antecedente en \citeA{martyniuk2025benchmarking}.

La última decisión de diseño que distingue esta propuesta de los enfoques habituales en la literatura es que los modelos no intentan predecir un único valor, sino cinco observables al mismo tiempo. La decisión de escoger varias en lugar de una se debe a que existen tipos de circuitos donde varios observables valen cero y no hay nada que predecir. Ampliando el número de observables a cinco, conseguimos que todas las familias de circutios aporten señales para el entrenamiento del modelo.


% =====================================================================================
% ⚠️ SECCION 1.3 «CONTRIBUCIONES» RETIRADA (sep-2026) POR DECISION DEL AUTOR.
%
% Se conserva comentada porque el material es reutilizable: las contribuciones 1 y 3
% estan enunciadas en §1.2.3, y la 2 y la 4 son los capitulos 8 y 7. Si se decide
% reponer la seccion, basta con descomentar el bloque de abajo y REVISARLO contra los
% resultados, que era su aviso original.
% =====================================================================================
% \section{Contribuciones}
% \label{sec:contribuciones}
% % Fuente: doc/memoria/materiales.md y doc/00_Justificaciones.md §4.
% % ⚠️ REVISAR AL CERRAR RESULTADOS.
% % ⚠️ Deliberadamente SIN CIFRAS de R²: las de materiales.md §... son del dataset v1 y la
% %    regla del proyecto prohibe citarlas. Reponer las del v2 cuando esten medidas.
%
% \begin{enumerate}
%     \item \textbf{Un conjunto de datos que no existía.} Veinte mil circuitos etiquetados con
%           su degradación exacta, generados sobre la calibración real de un procesador de IBM
%           a lo largo de $42$ días, y organizados en tres ejes de generalización separables
%           (tipo de circuito, tamaño y tiempo) que permiten medir por separado tres formas
%           distintas de fallar fuera de distribución.
%
%     \item \textbf{La demostración medida de que el objetivo natural del problema no es
%           predecible antes de ejecutar}, junto con su explicación: la degradación es el
%           producto de un factor que sí depende del circuito y del hardware por otro que
%           depende del algoritmo y que es, precisamente, la incógnita para la que haría falta
%           el ordenador cuántico. De ahí se sigue una reformulación del objetivo que sí es
%           aprendible, y que es la que sostiene el resto del trabajo.
%
%     \item \textbf{Una comparativa controlada de tres modelos que mide si la representación
%           en grafo aporta información real.} Ridge, Random Forest y una red neuronal de
%           grafos resuelven el mismo problema con el mismo protocolo y la misma partición; la
%           diferencia entre los dos primeros y el tercero es exactamente el valor de la
%           estructura que la agregación tabular descarta.
%
%     \item \textbf{La detección, medición y corrección de un error en las propias
%           etiquetas}, junto con la razón metodológica por la que la comprobación original no
%           lo detectó y el procedimiento con el que se validó la regeneración.
% \end{enumerate}

\section{Estructura de la memoria}
\label{sec:estructura}
% ⚠️ REESCRITA (sep-2026). La version anterior remitia a \ref{cap:fundamentos} y
%    \ref{cap:sota} como si fueran CAPITULOS, y desde la reestructuracion son las
%    SECCIONES 2.1 y 2.2: el texto imprimia «los capitulos 2.1 y 2.2». Corregido.
%
% 🔴 SI SE MUEVE, FUNDE O PARTE UN CAPITULO, HAY QUE REHACER ESTA LISTA. Es el unico
%    sitio de la memoria donde se enumeran los siete de corrido, asi que un desajuste
%    aqui se ve a la primera.
%
% Guion acordado: una lista, un parrafo corto por capitulo, y que el numero de capitulos
% se vea de un vistazo. Nada de prosa corrida.

La memoria se organiza en \textbf{siete capítulos} y dos anexos:

\begin{enumerate}
    \item \textbf{Introducción.} Presenta para qué sirve un ordenador cuántico y por qué el
          ruido limita hoy lo que puede resolver, delimita el problema concreto que se aborda y
          enuncia la propuesta.

    \item \textbf{Contexto y estado del arte.} Reúne lo imprescindible para leer el resto —del
          qubit a las puertas y al circuito como grafo acíclico dirigido— y revisa lo publicado,
          ordenado por cuándo actúa cada familia de métodos respecto a la ejecución. Cierra
          marcando el hueco que este trabajo ocupa.

    \item \textbf{Objetivos y metodología de trabajo.} Fija el objetivo general y los cinco
          específicos, la métrica que decide entre modelos y el proceso con el que se llega a
          ellos.

    \item \textbf{Planteamiento de la comparativa.} Es el capítulo de diseño. Establece qué se
          predice (los observables, el factor de supervivencia y el umbral de señal), de dónde
          salen los datos y cómo se parten en tres ejes de generalización, y con qué métricas se
          juzga.

    \item \textbf{Desarrollo de la comparativa.} Ataca los tres modelos uno a uno: qué datos ve
          cada uno, qué preprocesado es exclusivamente suyo y con qué arquitectura e
          hiperparámetros se quedó.

    \item \textbf{Discusión y análisis de resultados.} Presenta lo que salió, separando dos
          preguntas distintas: cuánto acierta cada modelo al predecir el factor, y cuánto mejora
          el observable una vez aplicado.

    \item \textbf{Conclusiones y trabajo futuro.} Enlaza cada objetivo con su conclusión, declara
          las limitaciones del trabajo y propone las líneas que las cierran.
\end{enumerate}

Los anexos recogen las herramientas empleadas y la declaración de uso de herramientas de
inteligencia artificial.


% --- Contexto y estado del arte (10–15 páginas en total, es el tope de la norma) ------
% El título anterior era «Fundamentos de computación cuántica». Ya no vale: el capítulo
% cubre ahora también los modelos de aprendizaje automático. De paso desaparece el
% Overfull \hbox de 3,35 pt que obligaba a partir el título a mano.
\chapter{Fundamentos}
\label{cap:fundamentos}

% =====================================================================================
% GUION. Capitulo de CONTEXTO: recoge conocimiento que YA EXISTE y se cita.
% Reestructurado en ago-2026 en dos bloques, 2.1 cuantico y 2.2 aprendizaje automatico.
%
% 🔴 LA FRONTERA CON EL CAPITULO 9, que es la regla que evita duplicar:
%
%      2.2  → QUE ES el metodo. Conocimiento de la literatura. Se CITA.
%      cap. 9 → POR QUE ese metodo aqui, como se configura, que se midio.
%               Decisiones propias. Se JUSTIFICA con nuestros datos.
%
%    Ejemplo: «que es la regularizacion L2» va en 2.2; «por que Ridge y no minimos
%    cuadrados, dada la colinealidad extrema que medimos» va en el 9. Si una frase
%    contiene una cifra nuestra, esta en el capitulo equivocado.
%
% 🔴 TONO, y es asimetrico a proposito. El tribunal es de INTELIGENCIA ARTIFICIAL:
%      2.1 → no saben cuantica. Claro, ameno, con intuicion antes que formalismo.
%            El tecnicismo entra, pero siempre despues de haberlo hecho comprensible.
%      2.2 → hablas su idioma. Denso y rapido en Ridge y Random Forest, que conocen
%            mejor que nadie. Detenerse SOLO en la red de grafos, que es lo que
%            probablemente no han visto aplicado a esto.
%
% ⚠️ Reglamento art. 11.4.a: «no es original el trabajo que contenga un exceso de
%    informacion de otras fuentes, AUNQUE ESTEN DEBIDAMENTE CITADAS». En un capitulo
%    de fundamentos es donde mas facil es caer. Defensa: incluir solo lo que los
%    capitulos 3, 6 y 9 van a usar de verdad, y que cada apartado apunte hacia adelante.
%
% ⚠️ Solo se numera hasta \subsection (secnumdepth = 2 en la clase book). Un
%    \subsubsection saldria SIN numero. Por eso el cuarto nivel va como parrafo con
%    entradilla en negrita, igual que en el resto de la memoria.
% =====================================================================================

% Parrafo de entrada del capitulo: que hace falta saber para leer el resto y por que
% este capitulo tiene dos mitades que parecen inconexas y no lo son.
% ⚠️ La norma prohibe dos encabezados consecutivos sin texto entre medias.


% =====================================================================================
\section{Fundamentos de computación cuántica}
\label{sec:fund_cuantica}
% =====================================================================================
% Fuente principal: Nielsen y Chuang (2010). Citar con paginas concretas:
%   \cite[pp.~13--16]{nielsen2010quantum} · \citeA[cap.~4]{nielsen2010quantum}
% UNA sola entrada en bibliografia.bib para todo el libro; la pagina va en la cita.

% Parrafo de entrada de la seccion.

\subsection{Del bit al qubit}
\label{sec:bit_qubit}
% LA PUERTA DE ENTRADA. Si el tribunal se pierde aqui, se pierde en todo el capitulo.
% Arrancar por lo que SI conocen: el bit, y de ahi la diferencia.
%
% Fuente: Nielsen y Chuang (2010), seccion 1.2, p. 13.
% ⚠️ VERIFICAR LA PAGINA del ultimo parrafo (el estado |+> y el ejemplo del atomo
%    iluminado): puede estar en la p. 14 y no en la 13. No citar p. 13 para eso sin
%    comprobarlo en el libro.
El bit es el concepto fundamental de la computación y la información clásicas. La
computación y la información cuánticas se construyen sobre un concepto análogo, el bit
cuántico o, abreviadamente, el \textbf{qubit} \cite[secc.~1.2, p.~13]{nielsen2010quantum}.

¿Qué es entonces un qubit? Igual que un bit clásico tiene un estado —o 0, o 1—, un qubit
también tiene un estado. Dos estados posibles para un qubit son $\ket{0}$ y $\ket{1}$, que,
como cabe suponer, se corresponden con los estados 0 y 1 de un bit clásico. La notación
$\ket{\ }$ se denomina \textit{notación de Dirac} y es la
notación estándar para los estados en mecánica cuántica. La diferencia entre bits y qubits
es que un qubit puede encontrarse en un estado distinto de $\ket{0}$ o $\ket{1}$: también es
posible formar combinaciones lineales de estados, llamadas habitualmente
\textbf{superposiciones}:
\begin{equation}
\label{eq:superposicion}
\ket{\psi} = \alpha \ket{0} + \beta \ket{1}.
\end{equation}
Los números $\alpha$ y $\beta$ son números complejos.

Esa misma información se puede escribir de forma más simple. Como el estado queda determinado
por sus dos amplitudes, basta con apilarlas para obtener una \textbf{representación
vectorial} del qubit, tomando $\ket{0}$ y $\ket{1}$ como los dos vectores de la base:
\begin{equation}
\label{eq:vector_estado}
\ket{0} \equiv \begin{pmatrix} 1 \\ 0 \end{pmatrix},
\qquad
\ket{1} \equiv \begin{pmatrix} 0 \\ 1 \end{pmatrix},
\qquad
\ket{\psi} = \alpha\ket{0} + \beta\ket{1} \equiv
\begin{pmatrix} \alpha \\ \beta \end{pmatrix}.
\end{equation}
Las dos notaciones son equivalentes y a lo largo del documento se emplean indistintamente:
la de Dirac resulta más cómoda para razonar, y la vectorial para calcular, porque convierte
las operaciones sobre qubits en productos de matrices —como se verá en el
apartado~\ref{sec:puertas}—.

Un bit se puede examinar para determinar si está en el estado 0 o en el 1. Los ordenadores
lo hacen constantemente, por ejemplo cuando recuperan el contenido de su memoria. Cabe recalcar
que un qubit no se pueda examinar para determinar su estado cuántico, es decir, los
valores de $\alpha$ y $\beta$. En su lugar, la mecánica cuántica nos dice que solo podemos
adquirir una información mucho más restringida sobre el estado cuántico. Al medir un qubit
obtenemos o bien el resultado $0$, con probabilidad $|\alpha|^2$, o bien el resultado 1, con
probabilidad $|\beta|^2$. Naturalmente, $|\alpha|^2 + |\beta|^2 = 1$, puesto que las
probabilidades deben sumar uno \cite[p.~13]{nielsen2010quantum}.

Que un qubit pueda estar en un estado de superposición choca con nuestra comprensión de
«sentido común» del mundo físico que nos rodea. Un bit clásico es como una moneda: o cara,
o cruz. En monedas imperfectas puede haber estados intermedios, como quedarse en equilibrio
sobre el canto, pero en el caso ideal se pueden descartar. Un qubit, por el contrario, puede
existir en un continuo de estados entre $\ket{0}$ y $\ket{1}$, hasta que se lo observa
\cite[p.~13]{nielsen2010quantum}.

Aunque los qubits parezcan «raros» o contraintuitivos, son algo físico y real, no una idea
abstracta. Se han comprobado en muchos experimentos y existen distintas formas de
construirlos en el mundo real:
\begin{itemize}
    \item un qubit puede ser la polarización de un fotón, esto es, dos formas distintas de
          vibrar la luz;
    \item puede ser la orientación del espín de un núcleo dentro de un campo magnético;
    \item puede ser dos estados de un electrón en un átomo: el estado «fundamental»
          ($\ket{0}$) y el estado «excitado» ($\ket{1}$).
\end{itemize}

En este último ejemplo, el átomo funciona como un pequeño sistema de dos niveles. Si se le
ilumina con luz de la energía adecuada, se puede hacer que el electrón pase de $\ket{0}$ a
$\ket{1}$, o de $\ket{1}$ a $\ket{0}$. Hasta aquí, todo es clásico: dos estados, uno u otro.

Pero lo interesante llega después. Si no se le da suficiente tiempo de luz, el electrón no
llega completamente a $\ket{1}$, sino que se queda en un estado intermedio, una superposición
entre $\ket{0}$ y $\ket{1}$. Ese estado intermedio es lo que se denomina $\ket{+}$.

% PENDIENTE en esta misma subseccion, aun sin redactar:
%   - Definir |+> explicitamente, si se decide dar la formula: (|0> + |1>)/raiz(2).
%   - El entrelazamiento como segunda propiedad que da la ventaja computacional.
%   - Por que n qubits necesitan 2^n amplitudes, y por que eso es a la vez el argumento
%     que hace interesante el campo y el que limita el tamaño de nuestro conjunto de
%     datos (capitulo~\ref{cap:rigor}).
%   - La contrapartida que enlaza con el capitulo~\ref{cap:ruido}: el estado cuantico es
%     fragil, y esa fragilidad es el problema que este trabajo intenta mitigar.

\subsection{La esfera de Bloch y la información que guarda un qubit}
\label{sec:bloch}
% Fuente: Nielsen y Chuang (2010), seccion 1.2.
% ⚠️ VERIFICAR LAS PAGINAS: la esfera de Bloch y la figura 1.3 estan en la seccion 1.2,
%    pero NO en la p. 13 —caen algo despues, hacia las pp. 14-16—. Se cita por seccion
%    hasta comprobar el numero exacto en el libro. No inventar la pagina.
Una imagen útil para pensar en los qubits es la siguiente representación geométrica. Como
$|\alpha|^2 + |\beta|^2 = 1$, la ecuación~\eqref{eq:superposicion} se puede reescribir como
\begin{equation}
\label{eq:bloch_gamma}
\ket{\psi} = e^{i\gamma}\left( \cos\frac{\theta}{2}\ket{0}
             + e^{i\varphi} \sin\frac{\theta}{2}\ket{1} \right),
\end{equation}
donde $\theta$, $\varphi$ y $\gamma$ son números reales. El factor $e^{i\gamma}$ que aparece
delante se puede ignorar, porque no tiene efectos observables, y por esa razón se puede
escribir
\begin{equation}
\label{eq:bloch}
\ket{\psi} = \cos\frac{\theta}{2}\ket{0} + e^{i\varphi} \sin\frac{\theta}{2}\ket{1}.
\end{equation}

Los números $\theta$ y $\varphi$ definen un punto sobre la esfera unidad tridimensional,
tal como se muestra en la Figura~\ref{fig:bloch}. Esta esfera se denomina habitualmente
\textbf{esfera de Bloch}, y nos permite visualizar el estado de un solo
qubit. Muchas de las operaciones sobre qubits individuales se describen 
dentro de esta representación. Sin embargo, para varios qubits presenta una limitación
porque no se conoce ninguna generalización sencilla en la
esfera de Bloch \cite[secc.~1.2]{nielsen2010quantum}.

% ⚠️ FORMATO EXIGIDO por las Instrucciones, apartado 1.3, PAGINA 6:
%      «Titulo en la parte superior de la tabla o figura.»
%      «Fuente de la tabla o figura en la parte inferior.»
%    Por eso el \caption va ANTES del dibujo. La plantilla oficial lo hace al reves en
%    su ejemplo; aqui manda la norma escrita.
%
% FIGURA PROPIA, dibujada en TikZ. El codigo esta en figuras/cap02/bloch_tikz.tex,
% con sus parametros comentados (radio, perspectiva, y los angulos theta y phi del
% estado que se dibuja).
%
% Al ser de elaboracion propia NO arrastra ninguna licencia, y ademas usa la
% tipografia del documento, asi que las etiquetas casan con el resto del texto.
%
% ALTERNATIVA: la version descargada de Wikimedia sigue en figuras/cap02/bloch.pdf.
% Para volver a ella, sustituir el \input por
%     \includegraphics[width=0.42\textwidth]{cap02/bloch}
% y devolver el pie a: «Fuente: Glosser.ca (2012), Bloch Sphere.svg, Wikimedia
% Commons. Licencia CC BY-SA 3.0.»
\begin{figure}[htbp]
    \centering
    \caption{Un bit clásico frente a un qubit: dos estados posibles frente
             a estados infinitos sobre la esfera de Bloch.}
    \label{fig:bloch}
    % =====================================================================================
% Esfera de Bloch — figura propia en TikZ.
%
% Se incluye desde capitulos/02_fundamentos.tex con \input. Es vectorial, usa la
% tipografia del documento y es de ELABORACION PROPIA: no arrastra ninguna licencia.
%
% GEOMETRIA. El plano ecuatorial se dibuja como una elipse alineada con los ejes de
% la pagina, de semiejes \R y \K*\R. Un punto del ecuador se parametriza con el angulo
% de pantalla t:            (\R cos t , \K \R sin t)
% Los ejes se colocan sobre esa misma elipse, asi que todo queda consistente:
%     \Ty  = 0°    -> el eje y, hacia la derecha
%     \Tx  = 235°  -> el eje x, hacia abajo a la izquierda
% El angulo azimutal phi se mide DESDE el eje x hacia el eje y, luego t = \Tx + phi.
%
% PARAMETROS AJUSTABLES: \R (radio), \K (aplastamiento de la perspectiva),
% \Theta y \Phi (posicion del estado dibujado).
% =====================================================================================
% ⚠️ \shorthandoff{<>} ES NECESARIO. `babel` con la opcion `spanish` convierte < y >
%    en caracteres ACTIVOS (los usa para las comillas << >>), y eso rompe la sintaxis
%    de puntas de flecha de TikZ: «>={Stealth}» y «[->]» fallan con
%    «Argument of \language@active@arg> has an extra }».
%    Va dentro de un grupo, asi que los catcodes se restauran solos al cerrar.
\begingroup
\shorthandoff{<>}%
\begin{tikzpicture}[>={Stealth[length=2.2mm]}, line width=0.5pt]

  \pgfmathsetmacro{\R}{2.2}      % radio de la esfera
  \pgfmathsetmacro{\K}{0.30}     % achatamiento del ecuador en perspectiva
  \pgfmathsetmacro{\Tx}{235}     % angulo de pantalla del eje x
  \pgfmathsetmacro{\Theta}{40}   % angulo polar del estado dibujado
  \pgfmathsetmacro{\Phi}{55}     % angulo azimutal del estado dibujado

  % --- coordenadas del estado |psi> ------------------------------------------------
  \pgfmathsetmacro{\Tp}{\Tx+\Phi}                       % angulo de pantalla del azimut
  \pgfmathsetmacro{\Rxy}{\R*sin(\Theta)}                % radio de la proyeccion ecuatorial
  \pgfmathsetmacro{\Px}{\Rxy*cos(\Tp)}
  \pgfmathsetmacro{\Pyeq}{\K*\Rxy*sin(\Tp)}             % altura de la proyeccion
  \pgfmathsetmacro{\Py}{\Pyeq+\R*cos(\Theta)}           % altura del propio estado

  % =====================================================================================
  % EL BIT CLASICO, a la izquierda, para la comparativa.
  % ⚠️ NO TOCA NADA DE LA ESFERA: se dibuja en su propia columna, a la izquierda del
  %    origen. La esfera sigue centrada en (0,0) exactamente como estaba.
  %
  % Los dos puntos van al MISMO azul que los polos de la esfera, y alineados con ellos.
  % Es deliberado: asi se lee de un vistazo que el bit clasico tiene DOS estados y que
  % el qubit tiene esos mismos dos... y toda la superficie que hay entre medias.
  % =====================================================================================
  \pgfmathsetmacro{\Xbit}{-2.60*\R}     % columna del bit clasico
  % ⚠️ Los cuatro puntos van al MISMO radio (\Rpunto) y las dos parejas a la MISMA
  %    altura, y=+\R y y=-\R: asi el 0 queda alineado con |0> y el 1 con |1>.
  \pgfmathsetmacro{\Rpunto}{3.2}        % radio de los puntos, en pt

  \fill[rojoclasico] (\Xbit,\R) circle (\Rpunto pt);
  \node[right=5pt] at (\Xbit,\R) {$0$};
  \fill[rojoclasico] (\Xbit,-\R) circle (\Rpunto pt);
  \node[right=5pt] at (\Xbit,-\R) {$1$};

  % rotulos de cada mitad
  \node[font=\bfseries] at (\Xbit,{-1.62*\R}) {Bit clásico};
  \node[font=\bfseries] at (0,{-1.62*\R}) {Qubit};

  % --- esfera y ecuador -------------------------------------------------------------
  \draw (0,0) circle (\R);
  \draw[dashed] (0,0) ellipse ({\R} and {\K*\R});

  % --- ejes -------------------------------------------------------------------------
  % Las puntas de flecha MUEREN sobre la propia esfera, sin sobresalir; las etiquetas
  % van fuera. El eje x es mas corto a proposito: no apunta hacia el borde del circulo
  % sino hacia el punto del ECUADOR que le corresponde, que en perspectiva cae dentro.
  % eje z: atraviesa la esfera de polo a polo, con punta arriba
  \draw (0,-\R) -- (0,0);
  \draw[->] (0,0) -- (0,\R);
  % El eje z lleva SU PROPIA etiqueta en negro, como x e y. Las dos etiquetas azules
  % de arriba y abajo no nombran el eje: nombran los dos ESTADOS de la base.
  \node[right=2pt] at ({0.02*\R},{0.86*\R}) {$\hat{\mathbf{z}}$};
  % Solo el estado: el eje ya lo nombra la etiqueta negra de arriba. Centradas sobre
  % el polo, que es lo que marcan.
  \node[above=2pt, azulunir] at (0,\R) {$\ket{0}$};
  \node[below=2pt, azulunir] at (0,-\R) {$\ket{1}$};
  % eje y: hacia la derecha, hasta el borde
  \draw[->] (0,0) -- (\R,0);
  \node[right=2pt] at (\R,0) {$\hat{\mathbf{y}}$};
  % eje x: hasta su punto del ecuador, abajo a la izquierda
  \draw[->] (0,0) -- ({\R*cos(\Tx)},{\K*\R*sin(\Tx)});
  \node[below left, inner sep=1pt] at ({\R*cos(\Tx)},{\K*\R*sin(\Tx)})
        {$\hat{\mathbf{x}}$};

  % --- los dos estados de la base computacional --------------------------------------
  % |0> y |1> son EXACTAMENTE los polos de la esfera. Marcarlos es correcto y ayuda:
  % deja ver de un vistazo que los estados clasicos son solo dos puntos de todos los
  % posibles. En azul institucional, el mismo de los titulos de capitulo.
  \fill[azulunir] (0,\R) circle (\Rpunto pt);
  \fill[azulunir] (0,-\R) circle (\Rpunto pt);

  % --- proyeccion del estado sobre el plano ecuatorial -------------------------------
  \draw[dotted] (0,0) -- (\Px,\Pyeq);
  \draw[dotted] (\Px,\Pyeq) -- (\Px,\Py);

  % --- el estado --------------------------------------------------------------------
  \draw[->] (0,0) -- (\Px,\Py);
  \fill (\Px,\Py) circle (1.5pt);
  \node[above right, inner sep=2pt] at (\Px,\Py) {$\ket{\psi}$};

  % --- angulo theta: entre el eje z y el estado --------------------------------------
  % El arco NO va de 90 a 90-theta: theta es el angulo real en 3D, y sobre el papel la
  % direccion del estado esta en atan2(\Py,\Px). Hay que arcar hasta ESE angulo para que
  % el arco muera justo sobre la flecha.
  \pgfmathsetmacro{\Apsi}{atan2(\Py,\Px)}
  \pgfmathsetmacro{\Rt}{0.46*\R}
  \draw (0,\Rt) arc[start angle=90, end angle=\Apsi, radius=\Rt];
  \pgfmathsetmacro{\Amid}{0.5*(90+\Apsi)}
  \node at ({1.24*\Rt*cos(\Amid)},{1.24*\Rt*sin(\Amid)}) {$\theta$};

  % --- angulo phi: en el plano ecuatorial, del eje x a la proyeccion -----------------
  % Se dibuja sobre la MISMA elipse del ecuador, a radio reducido, para que se lea como
  % un angulo tumbado en el plano y no como un angulo en el plano del papel.
  \pgfmathsetmacro{\Rp}{0.42*\R}
  \draw plot[domain=\Tx:\Tp, samples=40, smooth]
        ({\Rp*cos(\x)},{\K*\Rp*sin(\x)});
  \pgfmathsetmacro{\Pmid}{\Tx+0.5*\Phi}
  \node at ({1.55*\Rp*cos(\Pmid)},{1.55*\K*\Rp*sin(\Pmid)}) {$\phi$};

\end{tikzpicture}%
\endgroup


    {\footnotesize Fuente: elaboración propia.}
\end{figure}

¿Cuánta información representa un qubit? Teóricamente, existe un número infinito de
puntos sobre la esfera unidad. Sin embargo, esta
conclusión resulta engañosa, y el motivo es el comportamiento del qubit cuando se lo observa.
Conviene recordar que medir un qubit devuelve únicamente 0 o 1. Además, la medición cambia el
estado del qubit y lo colapsa, desde su superposición de $\ket{0}$ y $\ket{1}$, al estado
concreto. Por ejemplo, si la medición de
$\ket{+}$ da 0, el estado del qubit tras la medición será $\ket{0}$.

¿Por qué se produce este colapso? Nadie lo sabe: este comportamiento es sencillamente uno de
los postulados fundamentales de la mecánica cuántica. Lo relevante para nuestros propósitos es
que de una sola medición se obtiene un único bit de información sobre el estado del qubit.
Solo midiendo infinitos qubits preparados de forma idéntica
sería posible determinar $\alpha$ y $\beta$ para un qubit en el estado dado por la
ecuación~\eqref{eq:superposicion}.

Falta una pregunta todavía: ¿cuánta información representa un qubit si
no lo medimos? La pregunta es difícil porque un estado cuántico contiene información que no podemos medir directamente. 
Sin embargo, cuando un sistema cuántico cerrado evoluciona sin realizar ninguna medición, la naturaleza 
parece ``llevar la cuenta'' de todas las variables continuas que describen el estado, como $\alpha$ y $\beta$. 
En este sentido, incluso un solo qubit encierra una gran cantidad de información oculta, y dicha cantidad 
\textbf{crece de forma exponencial con el número de qubits}.
.

\subsection{Qubits múltiples}
\label{sec:multiples}
% Apartado BREVE y con poca matematica: su unico objetivo es que se vea el crecimiento
% 2^n antes de hablar de puertas, porque despues se usa constantemente.
% Fuente: Nielsen y Chuang (2010), seccion 1.2.
% ⚠️ VERIFICAR LA PAGINA antes de la entrega. Se cita por seccion.
Un solo qubit no basta para calcular nada interesante. Entonces, ¿cómo
se describe el estado conjunto de varios? La respuesta explica por sí sola por
qué un ordenador cuántico puede hacer cosas que uno clásico no.

Con \textbf{dos qubits} las combinaciones posibles de valores clásicos son cuatro: 00, 01,
10 y 11. A cada una le corresponde un estado de la base, y el estado conjunto es una
superposición de los cuatro, con una amplitud compleja por cada uno:
\begin{equation}
\label{eq:dos_qubits}
\ket{\psi} = \alpha_{00}\ket{00} + \alpha_{01}\ket{01}
           + \alpha_{10}\ket{10} + \alpha_{11}\ket{11}
\;\equiv\;
\begin{pmatrix} \alpha_{00} \\ \alpha_{01} \\ \alpha_{10} \\ \alpha_{11}\end{pmatrix}.
\end{equation}
La notación $\ket{00}$ es una abreviatura de $\ket{0}\otimes\ket{0}$, el
\textbf{producto tensorial} de los estados de cada qubit. La operación es más simple de lo
que sugiere el símbolo: consiste en multiplicar cada componente del primer vector por todas
las del segundo y apilar los resultados, de forma que dos vectores de dos componentes
producen uno de cuatro,
\begin{equation}
\label{eq:tensorial}
\begin{pmatrix} a \\ b \end{pmatrix} \otimes \begin{pmatrix} c \\ d \end{pmatrix}
=
\begin{pmatrix} a\,c \\ a\,d \\ b\,c \\ b\,d \end{pmatrix},
\qquad\text{y en particular}\qquad
\ket{0}\otimes\ket{0} =
\begin{pmatrix} 1 \\ 0 \end{pmatrix} \otimes \begin{pmatrix} 1 \\ 0 \end{pmatrix}
=
\begin{pmatrix} 1 \\ 0 \\ 0 \\ 0 \end{pmatrix} = \ket{00}.
\end{equation}
Es también la razón de que el número de amplitudes se multiplique por dos con cada qubit
que se añade: cada componente del estado anterior se desdobla en dos. Y la condición de
normalización se mantiene: la suma de los
módulos al cuadrado de las cuatro amplitudes vale uno, porque al medir se obtiene una de
las cuatro cadenas de bits con esas probabilidades.

Con \textbf{tres qubits} el razonamiento se repite. Las cadenas posibles pasan a ser ocho
—de 000 a 111—, y el estado necesita ocho amplitudes:
\begin{equation}
\label{eq:tres_qubits}
\ket{\psi} = \alpha_{000}\ket{000} + \alpha_{001}\ket{001} + \cdots + \alpha_{111}\ket{111}
\;\equiv\;
\begin{pmatrix} \alpha_{000} & \alpha_{001} & \cdots & \alpha_{111}\end{pmatrix}^{\!\intercal}.
\end{equation}

Y así sucesivamente:
\begin{itemize}
    \item con \textbf{cuatro qubits}, los estados de la base van de $\ket{0000}$ a
          $\ket{1111}$ y hacen falta \textbf{16 amplitudes};
    \item con \textbf{cinco qubits}, van de $\ket{00000}$ a $\ket{11111}$ y hacen falta
          \textbf{32 amplitudes};
    \item \dots
\end{itemize}

El patrón está claro: \textbf{$n$ qubits requieren $2^n$ amplitudes complejas}. Con 50
qubits son más de mil billones de números, y describir ese estado en un ordenador clásico
es imposible. Ahí está la promesa de la computación cuántica.

\subsection{Puertas cuánticas y circuitos}
\label{sec:puertas}
% Fuente de los dos primeros parrafos: Nielsen y Chuang (2010), seccion 1.3.
% ⚠️ VERIFICAR LA PAGINA antes de la entrega. Se cita por seccion porque no esta
%    comprobada; no inventarla.
%
% ⚠️ Del original se ha DESCARTADO la promesa «including a circuit which teleports
%    qubits!»: anuncia un contenido que esta memoria no va a desarrollar, y prometer
%    algo que luego no aparece es de las cosas que un tribunal si nota.
Los cambios que experimenta un estado cuántico se pueden describir con el lenguaje de la
computación cuántica. Del mismo modo que un ordenador clásico se construye a partir de un
circuito eléctrico formado por cables y puertas lógicas, un ordenador cuántico se construye
a partir de un circuito cuántico formado por cables y puertas cuánticas elementales,
encargadas de transportar y manipular la información cuántica
\cite[secc.~1.3]{nielsen2010quantum}.

Los circuitos de un ordenador clásico constan de cables y puertas lógicas. Los cables
sirven para transportar la información a lo largo del circuito, mientras que las puertas
lógicas realizan manipulaciones sobre esa información y la convierten de una forma en otra.

% La forma vectorial ya se introdujo en el apartado 2.1.1; aqui solo se recupera.
Escrita en forma vectorial —ecuación~\eqref{eq:vector_estado}—, la información de un qubit
es una columna de dos números complejos. Escrito así, el papel de una puerta cuántica se deduce solo: \textbf{una puerta es una
matriz que transforma el vector de estado en otro}. Si $U$ denota la puerta y $\ket{\psi}$
el estado de partida, el estado resultante es $\ket{\psi'} = U\ket{\psi}$, y aplicar varias
puertas en secuencia equivale a multiplicar sus matrices.

\paragraph{Un ejemplo.} La puerta $X$, análoga cuántica de la negación clásica, se define
como
\begin{equation}
\label{eq:puerta_x}
X = \begin{pmatrix} 0 & 1 \\ 1 & 0 \end{pmatrix},
\qquad\text{de modo que}\qquad
X\begin{pmatrix} \alpha \\ \beta \end{pmatrix} = \begin{pmatrix} \beta \\ \alpha \end{pmatrix}.
\end{equation}
Sobre los estados de la base actúa como cabría esperar de una negación, llevando $\ket{0}$
a $\ket{1}$ y $\ket{1}$ a $\ket{0}$. Sobre una superposición, en cambio, hace algo que no
tiene equivalente clásico: intercambia las dos amplitudes sin necesidad de medir, y por
tanto sin destruir la superposición.

% ⚠️ FORMATO: titulo arriba y fuente abajo (Instrucciones §1.3, p. 6).
\begin{table}[htbp]
    \centering
    \caption{Puertas cuánticas de un qubit más habituales.}
    \label{tab:puertas}
    \begin{tabular}{@{}llp{0.46\textwidth}@{}}
        \toprule
        \textbf{Puerta} & \textbf{Matriz} & \textbf{Efecto} \\
        \midrule
        $X$ & $\begin{pmatrix} 0 & 1 \\ 1 & 0\end{pmatrix}$
            & Intercambia $\ket{0}$ y $\ket{1}$. Es la negación cuántica. \\[1.1em]
        $Y$ & $\begin{pmatrix} 0 & -i \\ i & 0\end{pmatrix}$
            & Intercambia los estados e introduce además una fase. \\[1.1em]
        $Z$ & $\begin{pmatrix} 1 & 0 \\ 0 & -1\end{pmatrix}$
            & Deja $\ket{0}$ intacto y cambia el signo de $\ket{1}$. \\[1.1em]
        $H$ & $\dfrac{1}{\sqrt{2}}\begin{pmatrix} 1 & 1 \\ 1 & -1\end{pmatrix}$
            & Puerta de Hadamard: crea superposición a partir de un estado de la base. \\[1.1em]
        $S$ & $\begin{pmatrix} 1 & 0 \\ 0 & i\end{pmatrix}$
            & Aplica una fase de $\pi/2$ al estado $\ket{1}$. \\[1.1em]
        $T$ & $\begin{pmatrix} 1 & 0 \\ 0 & e^{i\pi/4}\end{pmatrix}$
            & Aplica una fase de $\pi/4$. \\[1.1em]
        $R_z(\theta)$ & $\begin{pmatrix} 1 & 0 \\ 0 & e^{i\theta}\end{pmatrix}$
            & Rotación de fase de ángulo arbitrario en torno al eje $z$. \\
        \bottomrule
    \end{tabular}

    {\footnotesize Fuente: elaboración propia.}
\end{table}

Como toda puerta de un qubit es unitaria, su efecto sobre la esfera de Bloch de la
Figura~\ref{fig:bloch} es siempre una \textbf{rotación}: el estado se desplaza sobre la
superficie. La Figura~\ref{fig:hadamard} es un ejemplo de como la puerta de
Hadamard, que puede descomponerse en un giro de $90^\circ$ en torno al eje $y$ seguido de
otro de $180^\circ$ en torno al eje $x$, actúa sobre el estado
$(\ket{0}+\ket{1})/\sqrt{2}$ y lo lleva de vuelta a $\ket{0}$.

% ⚠️ FORMATO: titulo arriba y fuente abajo (Instrucciones §1.3, p. 6).
% ⚠️ FUENTE: es un redibujo de una figura de Nielsen y Chuang, no una invencion propia.
%    La forma academica correcta es «elaboracion propia a partir de…», no «elaboracion
%    propia» a secas. Verificar el numero de figura del libro y añadirlo si se quiere.
% Codigo de color: azul = el qubit · naranja = la transformacion.
\begin{figure}[htbp]
    \centering
    \caption{Acción de la puerta de Hadamard sobre la esfera de Bloch.}
    \label{fig:hadamard}
    % =====================================================================================
% La puerta de Hadamard vista sobre la esfera de Bloch — figura propia en TikZ.
%
% Tres viñetas: el estado de partida (|0>+|1>)/raiz(2) sobre el eje x, la rotacion de
% 90 grados en torno a y que lo lleva al polo sur, y la rotacion de 180 grados en torno
% a x que lo sube al polo norte. Resultado: H|+> = |0>.
%
% CODIGO DE COLOR, deliberado y constante en toda la memoria:
%   azulunir   -> el QUBIT (vector de estado, punto y etiqueta)
%   acentofig  -> las TRANSFORMACIONES (giros, planos, flechas de operacion)
%
% ⚠️ \shorthandoff{<>} es obligatorio: `babel` en español hace activos < y >, y eso
%    rompe la sintaxis de puntas de flecha de TikZ.
% =====================================================================================
\begingroup
\shorthandoff{<>}%
\begin{tikzpicture}[>={Stealth[length=2mm]}, line width=0.5pt, scale=0.92]

  \pgfmathsetmacro{\R}{1.45}    % radio de cada esfera
  \pgfmathsetmacro{\K}{0.30}    % achatamiento del ecuador; ver nota al pie del fichero
  \pgfmathsetmacro{\Tx}{235}    % angulo de pantalla del eje x
  \pgfmathsetmacro{\D}{4.9}     % separacion entre viñetas

  % --- una esfera con sus tres ejes, reutilizable ------------------------------------
  \newcommand{\esferabase}{%
    \draw (0,0) circle (\R);
    \draw[dashed] (0,0) ellipse ({\R} and {\K*\R});
    \draw[->] (0,0) -- (0,\R);
    \node[right=1pt] at ({0.02*\R},{0.78*\R}) {\scriptsize $\hat{\mathbf{z}}$};
    \draw[->] (0,0) -- (\R,0) node[right=-1pt] {\scriptsize $\hat{\mathbf{y}}$};
    \draw[->] (0,0) -- ({\R*cos(\Tx)},{\K*\R*sin(\Tx)})
          node[below left=-3pt] {\scriptsize $\hat{\mathbf{x}}$};%
  }

  % =====================================================================================
  % 1) Estado de partida: (|0>+|1>)/raiz(2), sobre el eje x
  % =====================================================================================
  \begin{scope}[xshift=0cm]
    \esferabase
    \draw[azulunir, line width=1.1pt]
          (0,0) -- ({\R*cos(\Tx)},{\K*\R*sin(\Tx)});
    \fill[azulunir] ({\R*cos(\Tx)},{\K*\R*sin(\Tx)}) circle (2pt);
    \node[azulunir, anchor=east] at ({-1.06*\R},{-0.40*\R})
         {\small $(\ket{0}+\ket{1})/\sqrt{2}$};
    % giro de 90 grados en torno al eje y
    \draw[acentofig, ->, line width=0.7pt]
          ({0.30*\R},{-0.038*\R}) arc[start angle=-170, end angle=110,
          x radius={0.14*\R}, y radius={0.22*\R}];
  \end{scope}

  \draw[->, line width=0.8pt] ({0.5*\D-0.45},0) -- ({0.5*\D+0.45},0);

  % =====================================================================================
  % 2) Tras el giro: |1> en el polo sur. La rotacion de 180 grados en torno a x se
  %    representa con el plano y la flecha doble.
  % =====================================================================================
  \begin{scope}[xshift=\D cm]
    % el plano va PRIMERO: asi la esfera y el estado quedan por encima y se ven
    \draw[acentofig, fill=acentofig, fill opacity=0.06, line width=0.6pt]
          ({-1.78*\R},{-0.33*\R}) -- ({-0.90*\R},{0.31*\R})
       -- ({ 1.78*\R},{ 0.31*\R}) -- ({ 0.90*\R},{-0.33*\R}) -- cycle;
    \esferabase
    \draw[azulunir, line width=1.1pt] (0,0) -- (0,-\R);
    \fill[azulunir] (0,-\R) circle (2pt);
    \node[azulunir, anchor=north] at (0,{-1.08*\R}) {\small $\ket{1}$};
    % flecha doble: el estado atraviesa el plano de arriba abajo
    \draw[acentofig, <->, line width=0.9pt]
          ({0.58*\R},{-0.52*\R}) -- ({0.58*\R},{0.52*\R});
  \end{scope}

  \draw[->, line width=0.8pt] ({1.5*\D-0.45},0) -- ({1.5*\D+0.45},0);

  % =====================================================================================
  % 3) Resultado: |0> en el polo norte
  % =====================================================================================
  \pgfmathsetmacro{\DD}{2*\D}
    \begin{scope}[xshift=\DD cm]
    \esferabase
    \draw[azulunir, line width=1.1pt] (0,0) -- (0,\R);
    \fill[azulunir] (0,\R) circle (2pt);
    \node[azulunir, anchor=south east] at ({-0.04*\R},{1.06*\R}) {\small $\ket{0}$};
  \end{scope}

\end{tikzpicture}%
\endgroup

% =====================================================================================
% NOTA SOBRE \K, EL ACHATAMIENTO
%
% Determina desde que altura se mira la esfera, y con ello el angulo VISUAL entre los
% ejes x e y. Con el eje y horizontal, ese angulo no puede llegar a 90 grados: el eje x
% ha de apoyarse sobre la elipse del ecuador —porque el ecuador ES el plano xy— y la
% unica forma de que forme 90 grados con y seria que apuntase recto hacia abajo, encima
% del eje z. Lo medido:
%
%     K = 0,30  ->  157 grados   punta de x a 0,62·R del centro
%     K = 0,55  ->  142 grados   punta de x a 0,73·R del centro   <- el elegido
%     K = 0,75  ->  133 grados   punta de x a 0,84·R del centro
%     K = 1,00  ->  125 grados   punta de x a 1,00·R  (el ecuador se confunde con
%                                el contorno y la esfera deja de parecer una esfera)
%
% Se probo 0,55 y se descarto: el usuario prefiere la perspectiva original de 0,30.
% La tabla se deja aqui por si en algun momento se quiere volver a mover.
% =====================================================================================


    {\footnotesize Fuente: elaboración propia a partir de \citeNP{nielsen2010quantum}.}
\end{figure}

\paragraph{Ejemplo de puerta de dos qubits.} El estado conjunto de dos qubits necesita cuatro
amplitudes, una por cada combinación de la base $\ket{00}$, $\ket{01}$, $\ket{10}$ y
$\ket{11}$, de modo que una puerta que actúe sobre ambos es una matriz unitaria de
$4\times 4$. La más conocida es la \textbf{CNOT}, que niega el segundo qubit —el
\textit{objetivo}— solo cuando el primero —el \textit{control}— vale 1:
\begin{equation}
\label{eq:cnot}
\mathrm{CNOT} =
\begin{pmatrix}
1 & 0 & 0 & 0 \\
0 & 1 & 0 & 0 \\
0 & 0 & 0 & 1 \\
0 & 0 & 1 & 0
\end{pmatrix}.
\end{equation}

\paragraph{El circuito.} Un circuito cuántico es una secuencia de puertas aplicadas sobre
un conjunto de qubits, seguida de una medición. Se representa con una línea horizontal por
qubit —los «cables» del símil clásico— y las puertas colocadas sobre ellas en el orden en
que se ejecutan. Al circuito lo describen dos magnitudes: la \textbf{anchura},
que es el número de qubits, y la \textbf{profundidad}, que es el número de capas de puertas
que deben ejecutarse una tras otra.

La Figura~\ref{fig:circuito_ghz} muestra un ejemplo concreto. Los cinco cables horizontales
son los qubits; sobre el primero actúa una puerta de Hadamard y a continuación una cadena de
cuatro CNOT propaga el entrelazamiento al resto. El circuito tiene por tanto anchura cinco y
profundidad cinco, y prepara el estado
$(\ket{00000}+\ket{11111})/\sqrt{2}$, en el que los cinco qubits están entrelazados entre sí.

% ⚠️ FORMATO: titulo arriba y fuente abajo (Instrucciones §1.3, p. 6).
% Figura generada por el propio proyecto: TFM-Quantum/figures/00_anatomia/circuitos/.
% Se copia el PDF, no el PNG, porque es vectorial y no pixela.
%
% ⚠️ La imagen trae su PROPIO titulo incrustado («GHZ — Estado GHZ — el entrelazamiento
%    en su forma mas pura»), asi que se duplica con el \caption. Lo suyo es regenerarla
%    sin titulo desde src/viz/circuitos, o recortarlo. Se deja asi de momento.
\begin{figure}[htbp]
    \centering
    \caption{Tres ejemplos de circuito: preparación de un estado GHZ, un ansatz
             \textit{hardware-efficient} y una transformada cuántica de Fourier.}
    \label{fig:circuito_ghz}
    % Composicion compacta: GHZ y HEA en una fila, y la QFT —mas larga— debajo,
    % ocupando el ancho completo. Sin paquetes nuevos: dos minipages y un salto.
    % ⚠️ Las tres imagenes traen SU PROPIO titulo incrustado, que chocaba con el
    %    \caption y quedaba descolgado a la izquierda. Se recortan los 20 bp
    %    superiores de cada PDF —medido: el rotulo ocupa el 6,3 % del alto—.
    %    Lo limpio seria regenerarlas sin titulo desde TFM-Quantum/src/viz/circuitos.
    \includegraphics[trim=0 0 0 20bp, clip, width=0.47\textwidth]{cap02/circuito_ghz}\hfill
    \includegraphics[trim=0 0 0 20bp, clip, width=0.47\textwidth]{cap02/circuito_hea}

    \vspace{0.4em}
    \includegraphics[trim=0 0 0 20bp, clip, width=0.88\textwidth]{cap02/circuito_qft}

    {\footnotesize Fuente: elaboración propia.}
\end{figure}

\paragraph{La base nativa.} Un procesador real no implementa todas las puertas de la
Tabla~\ref{tab:puertas}, sino un conjunto reducido a partir del cual puede construirse
cualquier otra. Antes de ejecutarse, todo circuito se \textbf{traduce} a ese conjunto
mediante un proceso llamado \textit{transpilación}. En los procesadores IBM Heron la base
nativa es $\{cz, id, rz, sx, x\}$.

Esa traducción evidencia más problema que aborda este trabajo. Un
circuito teóricamente tiene diez puertas puede convertirse en cientos al expresarse en
la base nativa, y \textbf{cada puerta ejecutada es una oportunidad más de introducir
error}. Cuanto más largo es el circuito traducido, menos señal sobrevive hasta la medición.

% =====================================================================================
% PENDIENTE en esta subseccion:
%   - FIGURA de un circuito sencillo (dos qubits, H + CNOT + medicion). Seria muy util
%     para un tribunal que no sabe de cuantica: ahora mismo el circuito se describe con
%     palabras pero no se ve. Se puede dibujar en TikZ, como la esfera de Bloch, o con
%     el paquete `quantikz` si se decide añadirlo al preambulo.
%   - Verificar la pagina exacta de la cita de Nielsen y Chuang.
% =====================================================================================

\subsection{El circuito como grafo acíclico dirigido}
\label{sec:dag}
% ⚠️ ESTA SUBSECCION ERA LA BISAGRA hacia el 2.2: cerraba el bloque cuantico con la idea
%    de que, si un circuito es un grafo, predecirlo es regresion sobre grafos. Al subirla
%    aqui quedan DOS apartados por debajo (hardware y utilidad), asi que la transicion ya
%    no cae junto al 2.2. Hay que reponerla al final de «Para que sirve un ordenador
%    cuantico», o el 2.2 arrancara en seco.
Un circuito \textbf{es} un DAG: cada puerta es un nodo, las dependencias son las aristas.
Profundidad como camino más largo.

% ⚠️ FORMATO: titulo arriba y fuente abajo (Instrucciones §1.3, p. 6).
% Ambas figuras las genera figuras/cap02/gen_circuito_dag.py, que las produce con el
% propio Qiskit: el grafo NO es una ilustracion inventada, es la representacion interna
% que el transpilador manipula de verdad. El coloreado de los nodos —verde entradas,
% azul operaciones, rojo salidas— lo pone `dag_drawer` por si mismo.
\begin{figure}[htbp]
    \centering
    \caption{Un circuito y su grafo acíclico dirigido: cada puerta es un nodo y cada
             arista, la dependencia sobre un qubit.}
    \label{fig:dag}
    \begin{minipage}[c]{0.44\textwidth}
        \centering
        \includegraphics[width=\textwidth]{cap02/dag_circuito}
    \end{minipage}\hfill
    \begin{minipage}[c]{0.52\textwidth}
        \centering
        \includegraphics[width=0.92\textwidth]{cap02/dag_grafo}
    \end{minipage}

    {\footnotesize Fuente: elaboración propia.}
\end{figure}

\textbf{Y aquí es donde la física se convierte en un problema de aprendizaje automático:} si
un circuito es un grafo, predecir su comportamiento es un problema de regresión sobre
grafos. Eso es lo que motiva la sección siguiente.


% =====================================================================================
% =====================================================================================
% 2.1.6 — EL HARDWARE ACTUAL
% Sustituye al bloque de comentarios que habia bajo \subsection{El hardware actual}.
% Enfoque acordado: divulgativo («a modo curiosidad»), sin tecnicismos, y quedandose
% solo con los datos que el resto de la memoria necesita: la topologia limitada, la
% recalibracion diaria y el coste de simular. Las cifras van fechadas en 2026 y con cita.
% =====================================================================================

\subsection{El hardware actual}
\label{sec:hardware}

Conviene hacerse una idea de qué son físicamente las máquinas de las que habla esta
memoria, aunque solo sea para situar con una idea general el problema. Un procesador cuántico
superconductor como los de IBM es un chip del tamaño de una moneda que vive en el fondo de
un criostato: un cilindro de cobre dorado, de casi dos metros, cuya única función es
mantener el chip a unas centésimas de grado por encima del cero absoluto. Los qubits no se
programan escribiendo en una memoria, sino disparándoles pulsos de microondas por unos
cables coaxiales; cada puerta del circuito es un pulso. Y la medición no devuelve un
resultado, sino una estadística.

% Parrafo nuevo: explica QUE significa «repetir miles de veces y contar». Es el que
% justifica que el objetivo del TFM sea un valor esperado y no una salida determinista.
Nos detenemos en esto último. Al medir un qubit no
se lee el estado en el que estaba: el estado \textbf{se destruye} y en su lugar aparece un
cero o un uno, elegido al azar con la probabilidad que marcaba ese estado. Una única
ejecución, por tanto, no informa de nada; devuelve una cadena de ceros y unos que podría
haber salido de otra máquina completamente distinta. Lo que se hace en la práctica es
ejecutar el mismo circuito miles de veces seguidas —cada repetición se llama
\emph{shot}— y anotar cuántas veces sale cada resultado. De ese recuento se obtiene el
promedio que interesaba, y el circuito se ha ejecutado entero, desde cero, en cada una de
esas miles de repeticiones. La consecuencia se todo esto es que la precisión es un factor que se paga caro, ya que para reducir el error
en los resultados es necesario ejecutar más veces los circuitos, lo que aumenta coste de uso del procesador.

Merece de interés mencionar dos características del hardware que lo hacen distinto de un ordenador clásico:

\textbf{Primera: los qubits no se hablan todos con todos.} En un chip clásico cualquier
dato puede llegar a cualquier registro; en uno cuántico cada qubit está conectado solo a
dos o tres vecinos. Si un algoritmo necesita que interactúen dos qubits que no lo son, la
única salida es ir acercando el estado de uno hasta el otro con puertas intermedias, y cada
puerta añadida es ruido añadido.

\textbf{Segunda: el chip no es el mismo dos días seguidos.} Estas máquinas se recalibran a
diario, y cada calibración publica un boletín con la tasa de error de cada puerta y los
tiempos de vida de cada qubit. Esos números cambian y un qubit que ayer era de los
buenos hoy puede ser de los malos, y hay puertas que se estropean y no se reparan.
Esta es la razón de que el conjunto de datos del Capítulo~\ref{cap:dataset}
contenga la calibración real de un intervalo de días distintos en lugar de una única
fotografía del chip.

% ⚠️ FORMATO: titulo arriba y fuente abajo (Instrucciones §1.3, p. 6).
% Tabla simplificada a proposito: solo el catalogo de chips. El comentario sobre el
% cambio de estrategia (subir a 1121 y bajar a 156) se retiro por peticion expresa.
\begin{table}[htbp]
    \centering
    \caption{Familias de procesadores cuánticos de IBM y su número de qubits.}
    \label{tab:familias}
    \begin{tabular}{@{}lrr@{}}
        \toprule
        \textbf{Familia} & \textbf{Año} & \textbf{Qubits} \\
        \midrule
        Falcon    & 2019 & 27 \\
        Eagle     & 2021 & 127 \\
        Osprey    & 2022 & 433 \\
        Condor    & 2023 & 1\,121 \\
        Heron     & 2023 & 133 / 156 \\
        Nighthawk & 2025 & 120 \\
        \bottomrule
    \end{tabular}

    {\footnotesize Fuente: elaboración propia a partir de \citeNP{ibm2026hardware}.}
\end{table}

% Bloque nuevo: la equivalencia chip cuantico <-> ordenador clasico.
% Cifras calculadas como 2^n amplitudes x 16 bytes (complejo de doble precision).
% El 17,2 GB de n=15 es dato del propio proyecto: con ruido el coste es 4^n = 2^30.

Aunque 156 qubits parecen poca cosa habría que repasar \textbf{cuánta memoria haría falta en un ordenador normal
para imitar el chip}. Y la respuesta es que describir el estado de $n$ qubits exige guardar
$2^n$ números, de modo que \textbf{cada qubit que se añade duplica} el coste. 
\begin{itemize}
    \item Con 30 qubits hacen falta unos 17 GB, la memoria de un portátil actual bueno.
    \item  Con 40 qubits, unos 17 TB, que ya es un servidor grande.
    \item Con 50 qubits, unos 18 PB: un centro de datos entero dedicado a imitar un chip del tamaño de una moneda. 
    \item Y con los 156 qubits de un Heron sale un número de 48 cifras que ninguna máquina actual ni futura podrá almacenar.
\end{itemize}

Estas ejecuciones con ruido de chips cuánticos se pueden simular con ordenadores clásicos mediante paquetes de 
simuladores, como Qiskit Aer. Debido al gran cómputo que exigen una ejecución cuántica, los 
circuitos simulados en este trabajo llegan hasta 15 qubits.

Con todo esto hemos visto que el chip real maneja 156 qubits sin problema y el ordenador clásico no es capaz
ni de acercarse. Por tanto, no se puede calcular de antemano cuánto va a degradar
el ruido un resultado, porque ese cálculo es precisamente lo que resulta imposible. Lo
que sí puede intentarse es \textbf{aprenderlo} a partir de ejemplos.

\section{Fundamentos de los modelos de aprendizaje automático}
\label{sec:fund_ia}
% =====================================================================================
% SOLO TEORIA. Que es cada familia de modelos y como funciona. Ni una cifra nuestra,
% ni una decision de diseño: eso es del capitulo 9.
%
% ⚠️ REFERENCIAS PENDIENTES DE VERIFICAR antes de citarlas aqui (no estan en el .bib):
%      - Ridge original: Hoerl y Kennard (1970), Technometrics
%      - Random Forest: Breiman (2001), Machine Learning
%      - Atencion / Transformer: Vaswani y cols. (2017), NeurIPS
%      - Un manual general de referencia (Hastie/Tibshirani/Friedman o Bishop)
%    ➡️ Verificar DOI de cada una. Regla del proyecto: ninguna cita sin DOI comprobado.

\subsection{Dos formas de representar un circuito}
\label{sec:dos_representaciones}
% APERTURA DE 2.2. Cinco lineas, pero son las que convierten las tres subsecciones
% siguientes en una comparativa en marcha y no en tres fichas sueltas.
Del grafo del apartado anterior se puede partir de dos maneras. La primera es
\textbf{resumirlo en una fila de números} —cuántas puertas, de qué tipo, cuánto dura— y
aplicar un modelo tabular clásico. La segunda es \textbf{dárselo al modelo tal cual}, con su
estructura intacta.

Esa disyuntiva es la que estructura todo el trabajo, y las tres familias que siguen son las
candidatas: dos que exigen la tabla y una que acepta el grafo.

\subsection{Regresión lineal regularizada: Ridge}
\label{sec:teoria_ridge}
Mínimos cuadrados ordinarios y qué optimizan. El problema de la \textbf{colinealidad}: cuando
dos columnas dicen casi lo mismo, la solución se vuelve inestable y los coeficientes se
disparan.

La \textbf{penalización $L_2$} y el compromiso sesgo--varianza: se acepta algo de sesgo a
cambio de coeficientes acotados. Interpretación geométrica.

\subsection{Ensembles de árboles: Random Forest}
\label{sec:teoria_rf}
El árbol de decisión como partición recursiva del espacio de entrada. Su problema: varianza
alta.

\textbf{Bagging} y aleatorización de variables en cada división: por qué promediar muchos
árboles decorrelacionados reduce la varianza sin aumentar apenas el sesgo.

\textbf{Una propiedad teórica que importará después:} la predicción es siempre un promedio de
valores vistos en el entrenamiento, así que el modelo \textbf{no puede extrapolar} fuera del
rango aprendido. Aquí solo se enuncia; sus consecuencias se miden en el
capítulo~\ref{cap:comparativa}.

\subsection{Redes neuronales sobre grafos}
\label{sec:teoria_gnn}
% AQUI SI hay que detenerse: es lo que el tribunal probablemente no ha visto aplicado
% a este problema.
%
% 🔴 CORREGIDO (sep-2026). Este apartado decia lo CONTRARIO de la decision vigente:
%    planteaba la atencion como alternativa al paso de mensajes. La arquitectura del
%    trabajo es un MPNN CON paso de mensajes y nodo virtual BIDIRECCIONAL; el Graph
%    Transformer se descarto, y con medidas. Ver doc/00_Justificaciones.md §5.
%
% ⚠️ AQUI SOLO VA LA TEORIA (que es una GNN, que es el paso de mensajes, que es un nodo
%    virtual). El POR QUE de la eleccion —la matriz de estructura constante en el 99 %,
%    las distancias de 30-70 pasos— es una decision propia y medida: va en el
%    capitulo~\ref{cap:comparativa}, no aqui.
Qué es una red neuronal sobre grafos y qué la distingue de una red sobre vectores o
imágenes.

El \textbf{paso de mensajes}: cada nodo agrega información de sus vecinos, y hacen falta
tantas capas como distancia se quiera cubrir.

\textbf{El nodo virtual} como forma de obtener una representación de todo el grafo: un nodo
ficticio conectado a los demás, cuyo estado final resume la estructura entera.

% PENDIENTE de redactar: mencionar tambien la ALTERNATIVA basada en atencion, aunque sea
% brevemente, porque el capitulo 4 cita GTraQEM y el 9 explica por que se descarto. Sin
% saber que existe, ese descarte no se entiende.

\chapter{Estado del arte}
\label{cap:sota}

% =====================================================================================
% GUION. Reescrito por completo en sep-2026.
%
% ⚠️ EL BORRADOR ANTERIOR ESTABA LLENO DE MARCADORES «[cite: 1, 2]» que NO son citas de
%    LaTeX: venian de una exportacion y se imprimian tal cual en el PDF. Sustituidos por
%    \cite{} de verdad.
%
% 🔴 RESTRICCION DURA DE EXTENSION. Los capitulos 2 y 3 son el bloque «Contexto y estado
%    del arte» y la norma lo acota a 10-15 PAGINAS EN TOTAL. El capitulo 2 ocupa 12, asi
%    que este debe caber en TRES. De ahi las decisiones de economia:
%      - cada trabajo se despacha en dos o tres frases, no en un parrafo;
%      - los paradigmas alternativos (Q-Cluster, Liao 2025) van en UNA frase dentro de
%        §3.2, no en seccion propia; su sitio es el capitulo de limitaciones;
%      - §3.5 va con TABLA: ahorra dos paginas de prosa y es justo lo que un tribunal
%        busca en un trabajo de tipo 3.
%    Si el capitulo 2 adelgaza, lo primero que merece crecer aqui es §3.3.
%
% Fuentes: TFM-Quantum/doc/SoTA/comparative_analysis.md (corpus principal) y
%          TFM-Quantum/doc/SoTA_ampliada/fichas/paper_{1..5}.md (solo la idea principal
%          de cada uno, por decision del autor).
%
% ⚠️ Reglamento art. 11.4.a: «no es original el trabajo que contenga un exceso de
%    informacion de otras fuentes, AUNQUE ESTEN DEBIDAMENTE CITADAS». Defensa de este
%    capitulo: cada trabajo se trae por lo que decide algo en los capitulos 5 a 9, y se
%    dice explicitamente que decide.
%
% ⚠️ Excluidos por no citables: Wang/Zhengzhou, Stirbu, Al Farib.
% =====================================================================================

Este capítulo recorre lo que ya existe para corregir el efecto del ruido, ordenado por
\textbf{cuándo} actúa cada familia de métodos respecto a la ejecución del circuito. Se ha seleccionado
este criterio porque es el que separa a este trabajo de todo lo demás.

\section{La mitigación clásica y su precio}
\label{sec:qem_clasico}

La corrección de errores completa (QEC) exige cientos de qubits físicos por cada qubit
lógico, inviable para las máquinas actuales \cite{preskill2018nisq}. La
mitigación de errores (QEM) es la alternativa: no elimina el ruido, sino que
estima el resultado que se habría obtenido sin él \cite{cai2023quantum}.

Las técnicas de referencia son la \textbf{extrapolación a ruido cero} (ZNE), que trata de estimar cuál habría sido el resultado de un circuito si el ruido del hardware no estuviera presente \cite{temme2017error}; la \textbf{cancelación probabilística de errores} (PEC), que utiliza métodos estadísticos para compensar los errores introducidos durante la ejecución del circuito, aunque requiere un mayor número de ejecuciones \cite{cai2023quantum}; y la \textbf{mitigación de lectura} (M3), centrada en corregir los errores que aparecen al medir el estado final de los qubits \cite{nation2023suppressing}. 

Sus planteamientos son distintos, pero comparten los dos rasgos
que este trabajo toma como punto de partida: todas \textbf{consumen ejecuciones adicionales}
del procesador, el recurso más escaso, y todas actúan \textbf{después} de ejecutar.

\section{Modelos predictivos}
\label{sec:ml_qem}

Una forma de evitar el elevado coste de las técnicas tradicionales de mitigación consiste en sustituir las ejecuciones extra por un modelo entrenado.
\citeA{liao2024machine} demostraron que una red puede aprender la corrección directamente a
partir del valor ruidoso, y validaron este enfoque en hardware IBM real hasta un centenar de
qubits. De sus resultados, sale un dato que este trabajo aprovecha: entre los modelos tabulares,
\textbf{Random Forest es la mejor referencia}, y por eso entra en la comparativa del
Capítulo~\ref{cap:comparativa}.

Sobre esa idea pero en la línea de grafos, GTraQEM \cite{bao2025gtraqem} representa el circuito como
un DAG (\textit{Directed Acyclic Graph}) y lo procesa con un \textit{graph transformer} apoyado en un \textbf{nodo virtual}
conectado a todos los demás, que recoge el contexto global. Esta idea del nodo virtual es una de las características adoptadas en este trabajo, aunque no se emplea la misma arquitectura por los motivos expuestos en el Capítulo~\ref{cap:comparativa}. QEMFormer \cite{bao2025qemformer} se encuentra entre los métodos más avanzados y combina información general del circuito con detalles sobre su estructura.

\textbf{Aquí se encuentra una diferencia importante en la memoria con respecto a esos trabajos:} todos ellos toman el \textbf{valor
ruidoso medido como entrada}, de modo que hay que haber ejecutado el circuito para poder
corregirlo. Además, los resultados de \citeA{placidi2026deep} indican que esta información es la que más contribuye al rendimiento del modelo y sin ella el modelo reduce considerablemente su precision.

\section{Técnicas pre-ejecución}
\label{sec:preejecucion}

Existe una familia distinta que decide \textbf{antes} de gastar la máquina. La herramienta
oficial de IBM, mapomatic \cite{nation2023suppressing}, analiza distintas formas de asignar un mismo circuito al hardware disponible y selecciona la que previsiblemente producirá menos errores; es una
heurística analítica, sin aprendizaje. \citeA{hartnett2024learning} dan un paso más implementando técnicas de aprendizaje
automático con datos medidos en hardware real.

Un linaje reciente sí predice magnitudes sobre el grafo del circuito antes de ejecutarlo, y es
lo más cercano a esta propuesta. Entre estos trabajos, QuEst \cite{wang2025quest} estima la fiabilidad esperada de un circuito, mientras que \citeA{liu2026output} utilizan una red basada en grafos para predecir directamente el resultado que producirá.  
SQuASH \cite{martyniuk2025benchmarking} compara una red de grafos con Random Forest sobre los mismos circuitos, \textbf{precedente directo de esta comparativa}. Y \citeA{tudisco2025graph}, antes de
proponer la suya, demuestran que el número de qubits, la profundidad y el recuento de puertas no bastan para separar los casos: hay que meter el circuito entero en el modelo en vez
de resumirlo.

Ninguno de los cinco predice la degradación de un observable, y ninguno modela el paso del tiempo.

\section{El hardware cambia con el tiempo}
\label{sec:drift}

\citeA{hirasaki2023detection} midieron que las tasas de error no se degradan de forma progresiva,
sino a saltos bruscos que aparecen en cualquier momento, y que la recalibración diaria
no basta para capturarlos. \citeA{huo2026anchor} lo confirman por otro lado: el mismo
programa ejecutado en momentos distintos devuelve resultados diferentes. Esto 
justifica que el conjunto de datos de esta memoria tenga un eje temporal. 


\section{El hueco en la literatura}
\label{sec:gap}

% ⚠️ FORMATO: la plantilla oficial pone el \caption DEBAJO. Se sigue la plantilla.
% ⚠️ Va anclada con [t] y en \small A PROPOSITO: con [htbp] y cuerpo normal no cabia al
%    pie de su pagina, flotaba a la siguiente y arrastraba el ultimo item de la lista,
%    dejando dos tercios de pagina en blanco y sacando el capitulo de sus tres paginas.
% Las cuatro columnas son los cuatro diferenciadores enunciados debajo de la tabla.
% ⚠️ Si alguna fila cambia, cambiarla tambien en doc/SoTA/comparative_analysis.md.
\begin{table}[t]
    \centering
    \small
    \begin{tabular}{@{}lcccc@{}}
        \toprule
        \textbf{Trabajo} & \textbf{Pre-ejec.} & \textbf{Magnitud} & \textbf{Multi-salida} & \textbf{Deriva} \\
        \midrule
        ZNE, PEC, M3                            & no & sí & no & no \\
        \citeNP{liao2024machine}                & no & sí & no & parcial \\
        GTraQEM, QEMFormer                      & no & sí & sí & no \\
        mapomatic, \citeNP{hartnett2024learning}& sí & no & no & no \\
        \citeNP{wang2025quest}                  & sí & sí & no & no \\
        \citeNP{liu2026output}                  & sí & sí & no & no \\
        \citeNP{aktar2024graph}                 & sí & sí & no & no \\
        \citeNP{martyniuk2025benchmarking}      & sí & sí & no & no \\
        \citeNP{tudisco2025graph}               & sí & no & no & no \\
        \midrule
        \textbf{Este trabajo}                   & \textbf{sí} & \textbf{sí} & \textbf{sí} & \textbf{sí} \\
        \bottomrule
    \end{tabular}

    \caption{El estado del arte frente a los cuatro rasgos que definen esta propuesta.}
    \label{tab:hueco}

    {\footnotesize Fuente: elaboración propia.}
\end{table}

La Tabla~\ref{tab:hueco} resume el corpus frente a los cuatro rasgos que definen este trabajo.

Ninguna fila reúne las cuatro columnas: 

\begin{enumerate}
    \item \textbf{Pre-ejecución} (la entrada del modelo nunca incluye el resultado medido).
    \item  \textbf{Magnitud}, que es lo único que permite corregir y no solo descartar.
    \item  La \textbf{multi-salida} sobre varios observables.
    \item El \textbf{paso del tiempo como eje }, con días de calibración reales y disjuntos entre entrenamiento y prueba.
\end{enumerate}


% --- Objetivos y metodología de trabajo ---------------------------------------------
\chapter{Objetivos y metodología de trabajo}
\label{cap:objetivos}

% =====================================================================================
% GUION. Reescrito en sep-2026 siguiendo el modelo de memoria que aporto el autor.
% Exigido por «Instrucciones para la redaccion del TFE» §2.5. Los tres elementos son
% OBLIGATORIOS: (1) objetivo general, (2) objetivos especificos, (3) metodologia.
%
% ⚠️ Los objetivos DEBEN ser SMART; la norma lo llama «imprescindible». Por eso §4.1
%    enuncia el objetivo en un parrafo y DESPUES lo desglosa letra por letra: es la forma
%    en que el tribunal puede comprobar el criterio sin tener que deducirlo.
%
% 🔴 DECISION DEL AUTOR (sep-2026) SOBRE EL CRITERIO DE EXITO: la metrica que decide es el
%    R² sobre f, NO la mejora sobre no mitigar.
%    Razon: el diferenciador del trabajo es la pre-ejecucion, y «lo mitigado» necesita el
%    valor ruidoso, que solo existe DESPUES de ejecutar.
%    ⚠️ CONSECUENCIA PENDIENTE: el capitulo~\ref{cap:comparativa} dice hoy lo contrario
%       («un modelo que no bata a no mitigar no sirve, por bueno que sea su R²»).
%       HAY QUE CORREGIRLO ALLI para que no se contradigan. Aqui manda el R².
%
% ⚠️ PENDIENTE: la FECHA de §4.1 es un marcador. No se ha inventado ninguna: hay que
%    poner la fecha real de entrega en cuanto se conozca. Aparece DOS veces (S/M y T).
%
% ⚠️ El objetivo general sigue PENDIENTE DE CONFIRMAR CON EL DIRECTOR.
% =====================================================================================

Para este trabajo se plantea un objetivo general, del que derivan objetivos
específicos. Después de definir los objetivos explicaremos la metodología empleada
para cumplirlos.

\section{Objetivo general}
\label{sec:objetivo_general}

El objetivo de este trabajo es realizar un \textbf{análisis comparativo de tres
técnicas de aprendizaje automático que predicen, antes de ejecutar un circuito cuántico y únicamente a
partir de su estructura y de la calibración del procesador, la degradación que el
ruido causará en sus resultados}. En dicho análisis se obtiene un conjunto de métricas
estándar ($R^2$ y MAE), a través de las que se determina cuál de las técnicas ofrece la
mejor solución, así como si estas técnicas superan una determinada «línea base».

Esta definición cumple los criterios «SMART», ampliamente aceptados como buena práctica para la fijación de objetivos:

\begin{itemize}
    \item \textbf{S: \textit{Specific}} (específico): comparar tres técnicas de aprendizaje automático (regresión Ridge y Random Forest y tres variantes de GNN) para predecir el \textbf{factor de supervivencia} de los observables tras la degradación por ruido sin ejecutar el circuito. Se observa además si superan la línea base.
    \item \textbf{M: \textit{Measurable}} (medible): la comparativa se realiza a través de un
          conjunto de métricas estándar (MAE, RMSE y $R^2$), desglosadas por los tres ejes de
          generalización y por cada uno de los cincos observables. \textbf{La métrica de referencia que decide entre
          los modelos es el $R^2$ sobre el factor de supervivencia}, medido sobre el conjunto
          de prueba. Se toma esta y no la mejora sobre el resultado ya mitigado porque esta
          última exige conocer el valor ruidoso, que solo existe después de ejecutar, y el
          trabajo se define precisamente por no ejecutar.
    \item \textbf{A: \textit{Attainable}} (alcanzable): se apoya en un estado del arte
          desarrollado en el Capítulo~\ref{cap:sota}, en un conjunto de datos ya generado de
          veinte mil circuitos con calibración real, y en tres técnicas de modelado documentadas.
    \item \textbf{R: \textit{Relevant}} (relevante): tal y como se ha expuesto en el
          Capítulo~\ref{cap:intro}, poder anticipar la degradación permite decidir qué
          circuitos merece la pena ejecutar sin consumir tiempo de procesador cuántico, un recurso muy caro y escaso en la actualidad.
    \item \textbf{T: \textit{Time-Related}} (con un tiempo determinado): se establece como
          fecha fin el 9 de agosto de 2026.
\end{itemize}

\section{Objetivos específicos}
\label{sec:objetivos_especificos}

Los objetivos específicos de este trabajo son los siguientes:

\begin{enumerate}
    \item \textbf{Comprensión de los fundamentos de la computación cuántica}
    \item \textbf{Construcción del conjunto de datos}, con calibración real del procesador y
          tres ejes de generalización independientes (tipo de circuito, tamaño y tiempo).
    \item \textbf{Planteamiento y análisis de la variable objetivo}, esto es, cómo representar el factor de degradación y comprobar si realmente es predecible antes de ejecutar, midiendo el $R^2$ en validación sobre cada candidato.
    \item \textbf{Realización de la comparativa} entre los tres modelos bajo un mismo protocolo, y determinación de un posible modelo óptimo.
    \item \textbf{Cuantificación del aporte de la estructura} del circuito frente a un resumen tabular.
\end{enumerate}

% ⚠️ Cada uno de estos objetivos debe cerrarse con una conclusion en el capitulo 11.
%    La norma lo pide explicitamente: «cada objetivo del trabajo se enlazara con una
%    conclusion».

\section{Metodología de trabajo}
\label{sec:metodologia}

Se sigue una metodología \textbf{empírico-cuantitativa}. Tras haber planteado el problema y haber revisado la información que existe sobre este, se establece la hipótesis de la investigación, a la que puede darse la siguiente forma negativa, a modo de
hipótesis nula: «\textit{ninguno de los modelos de la comparativa es capaz de predecir la
degradación por ruido de un circuito cuántico antes de ejecutarlo mejor que los demás, ni de
superar la línea base establecida por el valor medio de la variable objetivo}». Este estudio
trata de probar la falsedad de dicha hipótesis.


El desarrollo se estructura sobre \textbf{CRISP-DM} \cite{wirth2000crispdm}, el modelo de
proceso de referencia en proyectos de minería de datos, que sigue siendo el marco más
empleado dos décadas después de su publicación \cite{schroer2021crispdm}. La
Tabla~\ref{tab:crispdm} recoge cómo se concreta cada una de sus seis fases en este trabajo.

% ⚠️ FORMATO: la plantilla oficial pone el \caption DEBAJO. Se sigue la plantilla.
% ⚠️ Anclada con [t] y en \small: con [htbp] y cuerpo normal desplaza el capitulo una
%    pagina entera.
\begin{table}[t]
    \centering
    \small
    \begin{tabular}{@{}p{0.26\textwidth}p{0.66\textwidth}@{}}
        \toprule
        \textbf{Fase de CRISP-DM} & \textbf{Concreción en este trabajo} \\
        \midrule
        Comprensión del negocio & Revisión del estado del arte, identificación del hueco y
        formulación de la hipótesis nula. \\
        Comprensión de los datos & Análisis exploratorio de la calibración real del procesador
        y de las etiquetas generadas; detección de las magnitudes que no son predecibles. \\
        Preparación de los datos & Generación del conjunto, definición de la partición
        temporal, agregación a variables tabulares y construcción de los grafos. \\
        Modelado & Implementación y ajuste de los tres modelos con el mismo protocolo,
        buscando los hiperparámetros contra el conjunto de validación. \\
        Evaluación & Métricas desglosadas por los tres ejes y por observable, y contraste
        frente a la línea base. \\
        Despliegue & Publicación del código y de los artefactos, y documentación de los
        límites de la solución. \\
        \bottomrule
    \end{tabular}

    \caption{Concreción de las fases de CRISP-DM en este trabajo.}
    \label{tab:crispdm}

    {\footnotesize Fuente: elaboración propia a partir de \citeNP{wirth2000crispdm}.}
\end{table}

A modo de listado, para alcanzar los objetivos propuestos se siguen los siguientes
\textbf{pasos}:

\begin{enumerate}
    \item \textbf{Revisión del estado del arte} y de los trabajos relacionados.
    \item \textbf{Determinación de los modelos} que forman parte de la comparativa, a partir
          del análisis del estado del arte.
    \item \textbf{Diseño y generación del conjunto de datos}, con la calibración real del
          procesador y los tres ejes de generalización (tipo, tamaño y tiempo).
    \item \textbf{Análisis exploratorio} de los datos y \textbf{validación de las etiquetas},
          realizado únicamente sobre el conjunto de entrenamiento.
    \item \textbf{Determinación de la variable objetivo}, descartando las magnitudes que se
          midan como no predecibles.
    \item \textbf{Preparación de los datos}: partición temporal, agregación a variables
          tabulares y construcción de los grafos.
    \item \textbf{Implementación} de los tres modelos en Python.
    \item \textbf{Entrenamiento y ajuste} de los modelos. La búsqueda de hiperparámetros se
          realiza contra un conjunto de validación \textbf{temporal}, formado por los últimos
          días de calibración del entrenamiento, y no por validación cruzada aleatoria: una división al azar mezclaría días entre ambos conjuntos, inflando los resultados y ocultando cómo rendirá el modelo ante futuros datos.
    \item \textbf{Obtención de las métricas} de los distintos modelos: $R^2$ y MAE.
          desglosadas por eje de generalización y por observable.
    \item \textbf{Validación de los resultados} y contraste frente a la línea base.
    \item \textbf{Determinación de un posible modelo óptimo}, estrictamente mejor que el
          resto, o identificación de otras casuísticas.
    \item \textbf{Elaboración de las conclusiones.}
\end{enumerate}


\subsection{Tipo de trabajo}
\label{sec:tipo_trabajo}
Comparativa de soluciones (tipo 3). Líneas de trabajo: aprendizaje automático (Ridge y Random
Forest) y aprendizaje profundo (la red neuronal de grafos).

\subsection{Herramientas empleadas}
\label{sec:herramientas}

La generación del conjunto de datos y la simulación del ruido se realizan con \textbf{Qiskit}
y su simulador Aer \cite{javadi2024qiskit}. Los modelos tabulares se implementan con
\textbf{scikit-learn} \cite{pedregosa2011scikit}, y la red neuronal de grafos con
\textbf{PyTorch} \cite{paszke2019pytorch} y \textbf{PyTorch Geometric}
\cite{fey2019pyg}. El seguimiento de los experimentos, con su semilla, la versión del código
y la huella del conjunto de datos, se lleva con \textbf{Weights \& Biases}
\cite{biewald2020wandb}.

El cómputo se reparte entre la máquina de desarrollo y \textbf{Google Cloud}
\cite{google2026cloud}. Esta última es una necesidad: las muestras de
mayor tamaño exigen más memoria de la que tiene la máquina local, y el entrenamiento de la
red de grafos requiere una unidad de procesamiento gráfico.

\subsubsection{Declaración de uso de herramientas de inteligencia artificial}
% El detalle NO se repite aqui: vive en el anexo, que es donde la norma espera encontrar
% una declaracion formal. Aqui solo la mencion y el reenvio.
El trabajo ha empleado herramientas de inteligencia artificial generativa en tareas concretas
de programación, generación de figuras, redacción y búsqueda bibliográfica. El detalle
completo, junto con las tareas que no se han delegado, se recoge en el Anexo~\ref{anexo:ia}.


% --- Planteamiento de la comparativa ------------------------------------------------
% ⚠️ REESTRUCTURADO (sep-2026), por decisión del autor. Dos cambios:
%      1. La VARIABLE OBJETIVO va DELANTE del planteamiento: el capítulo 6 usa el factor
%         de supervivencia en la EDA, en el preprocesado y en las métricas, así que
%         definirlo después dejaba al lector leyendo métricas de una magnitud sin
%         presentar.
%      2. El capítulo 6 absorbe datos, EDA, preprocesado, modelos, entrenamiento y
%         criterios de éxito: es el bloque «Planteamiento» completo (Instrucciones §2.6).
%         El capítulo 7 pasa a ser la EJECUCIÓN de esa comparativa.
%    🔴 El antiguo capítulo «Rigor metodológico» quedó FUERA de la memoria por decisión
%       expresa del autor. No reintroducirlo: su material (el bug de las etiquetas del v1,
%       el límite de determinismo del transpilador y las dos trampas de medición) sigue en
%       TFM-Quantum/doc/changes/ si algún día se recupera.
\chapter{La variable objetivo}
\label{cap:variable}

% =====================================================================================
% BLOQUE «Planteamiento de la comparativa» (1/2). Es LA MEDIDA sobre la que se compara:
% la norma del tipo 3 exige identificarla antes del desarrollo, y por eso va DELANTE del
% capitulo 6, que la usa en la EDA, en el preprocesado y en las metricas.
%
% CAPITULO DIFERENCIAL: es el resultado propio central del trabajo.
% APROBADO POR EL DIRECTOR (ago-2026) en su forma de hallazgo.
%
% 🔴 ORIGEN DE CADA CIFRA — todas son del dataset v2 y todas son citables:
%      §5.1  -> notebooks/01_eda/01b_los_observables.ipynb   (train, n = 13.344)
%      §5.2  -> notebooks/01_eda/01c_la_variable_objetivo.ipynb  (ajuste train, medida val)
%              + doc/changes/2026-08_cambio_decisiones.md §1, §2 y §3 (submuestras de 3.000)
%      §5.3  -> notebooks/01_eda/01d_el_umbral.ipynb        (SOLO train)
%      §5.4  -> notebooks/01_eda/01e_anatomia_de_f.ipynb    (umbral 0,002 ya fijado)
%
%    ⚠️ NO se ha traido NADA de doc/hallazgo_objetivo.md: ese documento es del v1, sus
%       cifras no son citables y ademas describe un experimento sobre OCHO observables y
%       un umbral de 0,05 que ya no existen.
%
% 🔴 DOS COSAS RETIRADAS POR DECISION EXPRESA DEL AUTOR (sep-2026). No reponerlas:
%
%      1. La EXPLICACION ARITMETICA de por que Δ no se aprende —que Δ = (1−f)·⟨O⟩exacto
%         es el producto de un factor aprendible por otro que no lo es—. Ocupaba un
%         apartado propio detras de §5.2.2. El material sigue en doc/hallazgo_objetivo.md
%         §2 si algun dia se recupera.
%
%      2. TODA mencion a las tres dispersiones (`std_Z`, `std_X`, `std_Y`) que el dataset
%         guarda y el modelo no usa: son irrelevantes para el trabajo y meterlas es
%         relleno. Por eso §5.1.1 ya no dice que el fichero guarde ocho observables, y
%         §5.5 habla solo del termino afin.
%
% 🔴 FRONTERA CON EL CAPITULO 6, que es la regla que evita duplicar:
%      cap. 5 -> QUE es el objetivo y POR QUE es ese. Propiedades medidas de la etiqueta.
%      cap. 6 -> QUE SE HACE con el: la cota sobre la prediccion (§6.5), la perdida
%                robusta (§6.4), el tronco compartido (§6.3) y las metricas (§6.5).
%    Por eso aqui las cotas se DEDUCEN y se COMPRUEBAN sobre las etiquetas, y la regla
%    que sale de ellas se remite a §\ref{sec:cota} en vez de repetirse.
%
% ⚠️ EXTENSION. Si hay que bajarla, el orden de recorte es:
%      1. la Figura 5.4 (f frente al tamaño): su contenido cabe en dos frases;
%      2. la Tabla 5.2 (las tres candidatas): la Figura 5.2 ya la dice entera.
%    NO recortar §5.3: es la unica justificacion documentada del umbral.
%
% ⚠️ Solo se numera hasta \subsection (secnumdepth = 2 en la clase book). No usar
%    \subsubsection: saldria sin numero.
% =====================================================================================

En este capítulo se reproduce el razonamiento que se siguió para construir la variable
objetivo, un proceso que supuso grandes desafíos. Se empieza introduciendo el concepto de los
observables y por qué se eligieron esos cinco. Después se listan las tres candidatas a variable
objetivo, se mide cuál de ellas es aprendible y se explica cómo se reconstruye la degradación
sin renunciar a predecir antes de ejecutar. A continuación se justifica el umbral por debajo
del cual una muestra deja de tener señal que aprender y se dedica un par de páginas a describir
la anatomía de la elección y las decisiones que se tomaron respecto a su física. El capítulo
cierra con las limitaciones de las decisiones tomadas.


% =====================================================================================
\section{Los observables}
\label{sec:observables}
% =====================================================================================
% ⚠️ DECISION DE REDACCION: la seccion declara su fuente teorica UNA SOLA VEZ, en el
%    parrafo de entrada, exactamente igual que hace §2.1. Nielsen y Chuang cubre el
%    postulado de la medida, el valor esperado de un operador y los operadores de Pauli;
%    los cinco observables concretos son eleccion propia y se justifican con datos, no
%    con la cita.

Salvo que se indique lo contrario, la teoría de esta sección sale del mismo libro de cabecera del Capítulo~\ref{cap:fundamentos}: \textit{Quantum Computation and Quantum Information}, de \citeA{nielsen2010quantum}.


\subsection{Definición y elección}
\label{sec:cinco_observables}

Un circuito acaba en un estado $\ket{\psi}$ que no se puede leer: al medir se obtiene una
cadena de ceros y unos y el estado colapsa (apartado~\ref{sec:bit_qubit}). La única forma de
extraer información útil es repetir la ejecución y quedarse con un promedio. Lo que se
promedia es un \textbf{observable}, una magnitud física medible representada por un operador
$O$, y la cantidad que interesa es su \textbf{valor esperado}:

\begin{equation}
\label{eq:valor_esperado}
\langle O \rangle = \bra{\psi} O \ket{\psi} = \boldsymbol{\psi}^\dagger O \boldsymbol{\psi} ,
\end{equation}

\noindent No es una probabilidad, sino \textbf{un número real}.

Los observables de este trabajo se construyen con los tres \textbf{operadores de Pauli}, $X$,
$Y$ y $Z$, que son las proyecciones elementales que admite un qubit, una por cada eje de la esfera
de Bloch (Figura~\ref{fig:pauli}). Los tres tienen los mismos dos autovalores, $+1$ y $-1$: sea
cual sea el estado y sea cual sea el eje, una medida individual solo puede devolver uno
de esos dos números. De ahí sale la propiedad siguiente: el
valor esperado es un promedio de $+1$ y $-1$ y, por tanto, vive en $[-1,1]$ por construcción.

% ⚠️ FORMATO: titulo arriba y fuente abajo (Instrucciones §1.3, p. 6).
% FIGURA PROPIA en TikZ. Codigo en figuras/cap05/pauli_tikz.tex, con la MISMA geometria
% que figuras/cap02/bloch_tikz.tex para que las dos se lean como la misma esfera.
\begin{figure}[htbp]
    \centering
    \caption{Los tres operadores de Pauli sobre la esfera de Bloch. Cada uno mide una
             dirección distinta del mismo estado, y los tres devuelven $+1$ o $-1$.}
    \label{fig:pauli}

    % =====================================================================================
% Los tres operadores de Pauli sobre la esfera de Bloch — figura propia en TikZ.
%
% Tres viñetas con la MISMA esfera. En cada una se resalta el eje que ese operador mide y
% se marcan sus dos autoestados, que son los dos unicos resultados posibles de la medida:
%     Z -> eje vertical,      |0> (+1) arriba,        |1> (-1) abajo
%     X -> eje de la izq.,    |+> (+1),               |-> (-1)
%     Y -> eje de la derecha, |+i> (+1),              |-i> (-1)
% El mensaje es que medir X, Y o Z no es medir tres cosas, sino preguntarle al MISMO
% estado por tres direcciones distintas, y que las tres respuestas valen +1 o -1. De ahi
% que el valor esperado viva en [-1, 1], que es lo que gobierna el resto del capitulo 5.
%
% GEOMETRIA: la misma que en figuras/cap02/bloch_tikz.tex, para que las dos figuras se
% lean como la misma esfera. El ecuador es una elipse de semiejes \R y \K*\R, y un punto
% del ecuador se parametriza con el angulo de pantalla t: (\R cos t, \K \R sin t).
%     \Tx = 235° -> eje x, abajo a la izquierda      \Ty = 0° -> eje y, a la derecha
%
% CODIGO DE COLOR, el mismo que en el resto de la memoria:
%   azulunir   -> el QUBIT y sus estados
%   acentofig  -> lo que se le HACE: aqui, el eje que se esta midiendo
%   gris       -> los otros dos ejes, que en esa viñeta no se miden
%
% ⚠️ \shorthandoff{<>} es obligatorio: `babel` en español hace activos < y >, y eso rompe
%    la sintaxis de puntas de flecha de TikZ.
% =====================================================================================
\begingroup
\shorthandoff{<>}%
\begin{tikzpicture}[>={Stealth[length=1.9mm]}, line width=0.5pt, font=\footnotesize]

  \pgfmathsetmacro{\R}{1.25}    % radio de cada esfera
  \pgfmathsetmacro{\K}{0.30}    % achatamiento del ecuador en perspectiva
  \pgfmathsetmacro{\Tx}{235}    % angulo de pantalla del eje x
  \pgfmathsetmacro{\D}{4.45}    % separacion entre viñetas
  \pgfmathsetmacro{\Rp}{2.6}    % radio de los puntos, en pt
  \pgfmathsetmacro{\DD}{2*\D}   % desplazamiento de la tercera viñeta

  % --- coordenadas de los cuatro extremos ecuatoriales -------------------------------
  \pgfmathsetmacro{\Xa}{\R*cos(\Tx)}     \pgfmathsetmacro{\Ya}{\K*\R*sin(\Tx)}
  \pgfmathsetmacro{\Xb}{-\Xa}            \pgfmathsetmacro{\Yb}{-\Ya}

  % --- una esfera con sus tres ejes en gris, reutilizable ----------------------------
  \newcommand{\esferabase}{%
    \draw[gray!55] (0,0) circle (\R);
    \draw[gray!55, dashed] (0,0) ellipse ({\R} and {\K*\R});
    \draw[gray!55] (0,-\R) -- (0,\R);
    \draw[gray!55] (\Xa,\Ya) -- (\Xb,\Yb);
    \draw[gray!55] (-\R,0) -- (\R,0);%
  }

  % ===================================================================================
  % 1) Z — el eje vertical
  % ===================================================================================
  \begin{scope}[xshift=0cm]
    \esferabase
    \draw[acentofig, line width=1.2pt] (0,-\R) -- (0,\R);
    \fill[azulunir] (0,\R) circle (\Rp pt);
    \fill[azulunir] (0,-\R) circle (\Rp pt);
    \node[above=1pt, azulunir] at (0,\R)  {$\ket{0}$};
    \node[below=1pt, azulunir] at (0,-\R) {$\ket{1}$};
    \node[right=3pt, acentofig] at (0,{0.62*\R})  {$+1$};
    \node[right=3pt, acentofig] at (0,{-0.62*\R}) {$-1$};
    \node[font=\small] at (0,{-1.62*\R}) {$\langle Z \rangle$};
  \end{scope}

  % ===================================================================================
  % 2) X — el eje que en perspectiva sale hacia el lector
  % ===================================================================================
  \begin{scope}[xshift=\D cm]
    \esferabase
    \draw[acentofig, line width=1.2pt] (\Xa,\Ya) -- (\Xb,\Yb);
    \fill[azulunir] (\Xa,\Ya) circle (\Rp pt);
    \fill[azulunir] (\Xb,\Yb) circle (\Rp pt);
    \node[below left=-1pt,  azulunir] at (\Xa,\Ya) {$\ket{+}$};
    \node[above right=-1pt, azulunir] at (\Xb,\Yb) {$\ket{-}$};
    % Los autovalores se desplazan PERPENDICULARES al eje (hacia arriba y a la
    % izquierda), porque puestos encima del propio eje se pisaban con los kets.
    \node[acentofig] at ({0.62*\Xa-0.10},{0.62*\Ya+0.25}) {$+1$};
    \node[acentofig] at ({0.62*\Xb-0.10},{0.62*\Yb+0.25}) {$-1$};
    \node[font=\small] at (0,{-1.62*\R}) {$\langle X \rangle$};
  \end{scope}

  % ===================================================================================
  % 3) Y — el eje horizontal
  % ===================================================================================
  \begin{scope}[xshift=\DD cm]
    \esferabase
    \draw[acentofig, line width=1.2pt] (-\R,0) -- (\R,0);
    \fill[azulunir] (\R,0) circle (\Rp pt);
    \fill[azulunir] (-\R,0) circle (\Rp pt);
    \node[right=2pt, azulunir] at (\R,0)  {$\ket{{+}i}$};
    \node[left=2pt,  azulunir] at (-\R,0) {$\ket{{-}i}$};
    % Por encima de la elipse del ecuador: pegados al eje quedaban sobre la linea
    % de puntos y se leian mal.
    \node[acentofig] at ({0.58*\R},0.44)  {$+1$};
    \node[acentofig] at ({-0.58*\R},0.44) {$-1$};
    \node[font=\small] at (0,{-1.62*\R}) {$\langle Y \rangle$};
  \end{scope}

\end{tikzpicture}%
\endgroup


    {\footnotesize Fuente: elaboración propia.}
\end{figure}

Con esas tres piezas se puede construir cualquier observable de Pauli. Los elegidos para este trabajo son los cinco que se muestran en la Tabla~\ref{tab:observables}: las tres bases de Pauli promediadas sobre todos
los qubits, y dos correladores, uno global y otro local. Con estos cinco cubrimos lo básico: el fallo en cúbits individuales y los errores que surgen al relacionarlos.

% ⚠️ FORMATO: titulo arriba y fuente abajo (Instrucciones §1.3, p. 6).
% Definiciones tomadas de src/quantum_gen.py::_build_observables y ::_reduce. El orden es
% el de src/config.py::OBJETIVOS y NO debe cambiarse aqui sin cambiarlo alli.
\begin{table}[htbp]
    \centering
    \caption{Los cinco observables objetivo y qué mide cada uno.}
    \label{tab:observables}
    {\footnotesize
    \begin{tabular}{@{}llp{0.44\textwidth}@{}}
        \toprule
        \textbf{Observable} & \textbf{Definición} & \textbf{Qué mide} \\
        \midrule
        Media de $Z$ & $\frac{1}{n}\sum_{i} \langle Z_i \rangle$
            & La magnetización media en la base de medida: cuánto se inclinan los qubits
              hacia $\ket{0}$ o hacia $\ket{1}$. \\
        Media de $X$ & $\frac{1}{n}\sum_{i} \langle X_i \rangle$
            & Lo mismo en la base de las superposiciones. Es independiente de Z.\\
        Media de $Y$ & $\frac{1}{n}\sum_{i} \langle Y_i \rangle$
            & Con este se consigue las tres direcciones independientes. \\
        Paridad & $\langle Z_0 Z_1 \cdots Z_{n-1} \rangle$
            & Magnitud \textbf{global}: un solo operador sobre todos los qubits a la vez. \\
        Correlador de vecinos & $\frac{1}{n-1}\sum_{i} \langle Z_i Z_{i+1} \rangle$
            & Correlación \textbf{local}. \\
        \bottomrule
    \end{tabular}}

    {\footnotesize Fuente: elaboración propia.}
\end{table}

\subsection{Impacto del ruido a un valor esperado}
\label{sec:ruido_observable}

Los cinco observables están acotados en todo el intervalo $[-1,1]$, pero la mayor parte se acumula en un intervalo más reducido $[-0.5,\, 0.5]$, como se puede observar en el panel izquierdo de la Figura~\ref{fig:observables_ruido}. Además, ese mismo panel sugiere que los valores exactos tras la ejecución se reducen, ya que los rangos usuales de los valores ruidosos son menores. Sin embargo, como se muestra en el panel derecho de la Figura~\ref{fig:observables_ruido}, el ruido afecta de dos maneras:
\begin{itemize}
    \item \textbf{Fenómeno de encogimiento}: se produce en gran parte de los circuitos (entre el $86{,}6$\% y el $92{,}4$\% de las veces).
    \item \textbf{Fenómeno de expansión}: se produce en menor número de circuitos (entre el $7{,}6$\% y el $13{,}4$\%).
\end{itemize}

% ⚠️ FORMATO: titulo arriba y fuente abajo (Instrucciones §1.3, p. 6).
% Figura propia. La genera observable.dibujar_rangos_y_efecto() en el cuaderno
% notebooks/01_eda/01b, y se copia desde figures/eda/observables/02_rangos_y_efecto.pdf.
% ⚠️ Va SIN titulo dentro de la imagen a proposito: el titulo es el \caption, que es lo que
%    pide la norma. La funcion `dibujar_rangos` original, que si lleva titulo, se conserva
%    en src/eda/observable.py para el uso interno del cuaderno.
% ⚠️ El panel izquierdo NO dibuja los atipicos (`showfliers=False`), asi que su eje no
%    llega a +-1 aunque el observable si. No escribir a partir de el ninguna afirmacion
%    sobre el rango COMPLETO: para eso esta la tabla de `observable.rangos()`.
\begin{figure}[htbp]
    \centering
    \caption{Los cinco observables frente al ruido. Izquierda: dónde vive cada uno, con y sin
             ruido. Derecha: en cuántos circuitos el ruido lo encogió y en cuántos lo agrandó.}
    \label{fig:observables_ruido}

    \includegraphics[width=\textwidth]{cap05/obs_rangos_y_efecto}

    {\footnotesize Fuente: elaboración propia.}
\end{figure}

Entonces el ruido puede amplificar y no es una anomalía de la simulación.

% =====================================================================================
\section{El factor de supervivencia}
\label{sec:factor}
% =====================================================================================

\subsection{Candidatas a variable objetivo y elección}
\label{sec:candidatas}

La pregunta del trabajo, \emph{¿cuánto degrada el ruido?}, tuvo 3 opciones a poder responderlo:

\begin{enumerate}
    \item \textbf{La degradación con signo}, que era el objetivo de inicial del proyecto:
          \begin{equation}
          \label{eq:delta}
          \Delta = \langle O \rangle_{\text{exacto}} - \langle O \rangle_{\text{ruidoso}} .
          \end{equation}
          Es lo que se quiere corregir y pedecirla permitiría sumarla al resultado medido.

    \item \textbf{La magnitud de la degradación}, $|\Delta|$, que renuncia a la dirección del
          error y no sirve para corregir.

    \item \textbf{El factor de supervivencia}, la fracción de señal que queda tras la
          ejecución:
          \begin{equation}
          \label{eq:f}
          f = \frac{\langle O \rangle_{\text{ruidoso}}}{\langle O \rangle_{\text{exacto}}} .
          \end{equation}
          \begin{itemize}
            \item $f = 1$, si el cirucito salió intacto.
            \item $f < 1$, si el ruido redujo.
            \item $f > 1$, si el ruido amplificó.
            \item $f = 0$, destruyó la señal completamente.
          \end{itemize}
\end{enumerate}

La elección se midió. Se ajustó el mismo modelo, con las mismas variables y la misma partición, cambiando únicamente el objetivo, y se evaluó en validación:
mismos tipos de circuito y mismo rango de tamaños que el entrenamiento. Se usaron los dos
modelos tabulares de la comparativa y se repitió todo con tres umbrales de señal distintos (la explicación de estos umbrales se realiza en el Capítulo~\ref{sec:umbral}), para comprobar que la conclusión no dependía de esa constante
(Tabla~\ref{tab:candidatas} y Figura~\ref{fig:objetivo_r2}).

\begin{table}[htbp]
    \centering
    \caption{Coeficiente de determinación en validación de las tres cadidatas (umbral fijo $0.01$). Cada celda
             es el recorrido sobre los cinco observables.}
    \label{tab:candidatas}
    {\footnotesize
    \begin{tabular}{@{}lcc@{}}
        \toprule
        \textbf{Objetivo} & \textbf{Ridge} & \textbf{Random Forest} \\
        \midrule
        $\Delta$ con signo & $-0{,}019$ \dots $+0{,}005$ & $-0{,}175$ \dots $-0{,}021$ \\
        $|\Delta|$         & $+0{,}260$ \dots $+0{,}366$ & $+0{,}167$ \dots $+0{,}303$ \\
        \textbf{Factor de supervivencia} & $\bm{+0{,}657}$ \dots $\bm{+0{,}782}$
                                         & $\bm{+0{,}652}$ \dots $\bm{+0{,}804}$ \\
        \bottomrule
    \end{tabular}}

    {\footnotesize Fuente: elaboración propia. Umbral de señal $0{,}01$.}
\end{table}

% ⚠️ FORMATO: titulo arriba y fuente abajo (Instrucciones §1.3, p. 6).
% Figura propia. La genera variable_objetivo.dibujar_comparativa_memoria() en el cuaderno
% notebooks/01_eda/01c, y se copia desde
% TFM-Quantum/figures/eda/objetivo/02_variable_objetivo_memoria.pdf.
% ⚠️ Va SIN titulo dentro de la imagen: lo pone el \caption, que es lo que pide la norma.
%    Los rotulos de cada panel SI se quedan, porque no titulan la figura.
% ⚠️ Su `figsize` esta calculado para la caja de 15 cm: a \textwidth se imprime a escala
%    ~0,6, y por eso sus cuerpos de letra son de 11-12 pt en origen. Agrandar la figura sin
%    subir las fuentes deja la letra ilegible en el PDF.
\begin{figure}[htbp]
    \centering
    \caption{Las tres candidatas del objetivo, por observable y con tres umbrales de señal
             distintos. Random Forest, coeficiente de determinación en validación.}
    \label{fig:objetivo_r2}

    \includegraphics[width=\textwidth]{cap05/objetivo_r2}

    {\footnotesize Fuente: elaboración propia.}
\end{figure}

El resultado no admite matices de lectura:

\begin{itemize}
    \item \textbf{La degradación con signo no se aprende.} Su coeficiente de determinación es
          nulo o negativo en los cinco observables, con los dos modelos y con los tres
          umbrales.
    \item \textbf{La magnitud sí se aprende, pero poco.} Renunciar al signo recupera señal predecible. Pero un objetivo sin signo no corrige nada.
    \item \textbf{El factor de supervivencia es el único aprendible de verdad}, y duplica a $|\Delta|$ en los dos modelos y en los tres umbrales.
\end{itemize}

\subsection{La reformulación}
\label{sec:reformulacion}

De las tres candidatas, la única que el modelo puede aprender es el factor de supervivencia, y
esa es la que se escoge como objetivo de entrenamiento. La degradación se
reconstruye después, a partir del factor predicho y del valor medido tras ejecución:
\begin{equation}
\label{eq:reconstruccion}
\hat\Delta = \langle O \rangle_{\text{ruidoso}}\cdot\frac{1 - \hat f}{\hat f} .
\end{equation}

% ⚠️ EL SIGNO ES UNA SUMA, y es un error facil de arrastrar. Delta esta definido en el
%    dataset como Delta = exacto − ruidoso, luego exacto = ruidoso + Delta. Con una resta,
%    un modelo PERFECTO duplicaria el error en vez de anularlo. Corregido en sep-2026;
%    hasta entonces cinco documentos del repositorio ponian la resta.

\textit{Una acalaración}: para reconstruir la degradación, en gran parte de la literatura \citeA{placidi2026deep} recurren al valor ruidoso $\langle O \rangle_{\text{ruidoso}}$ (obtenido tras ejecución en \textbf{hardware cuántico real} y costoso pero acanzable) y del valor $\langle O \rangle_{\text{exacto}}$ (obtenido tras ejecución en \textbf{hardware clásico} y costoso pero con alcance limitado, explicado en el Capítulo~\ref{sec:hardware})



\begin{center} 
      \doublebox{
      \begin{minipage}{0.85\linewidth} \vspace{0.5em}
             \centering 
             \textit{Sin embargo, en nuestra propuesta nunca haría falta calcular el valor $\langle O \rangle_{\text{exacto}}$ para poder mitigar el ruido, donde existen actualmente limitaciones a partir de un cierto número de qubits.} \vspace{0.5em} \end{minipage}} 
\end{center}


A partir de la igualdad \ref{eq:reconstruccion}, obtenemos la reconstrucción del valor mitigado así:

\begin{equation}
\label{eq:mitigado}
\langle O \rangle_{\text{mitigado}}
  = \langle O \rangle_{\text{ruidoso}} + \hat\Delta
  = \frac{\langle O \rangle_{\text{ruidoso}}}{\hat f} .
\end{equation}

Conviene profundizar en la idea propuesta, ya que este plantemaiento no rompe el \textbf{paradigma pre-ejecución}. Las razones son las siguientes:

\begin{itemize}
    \item \textbf{La entrada del modelo} es el circuito y la telemetría del chip, y nada más. El modelo \textbf{nunca ve} el resultado de la ejecución, ni al entrenar ni al predecir.
    \item \textbf{El valor medido entra solo tras la predicción} para mitigar el ruido si corresponde la decisión.
    \item \textbf{Y la predicción sirve sin ejecutar nada.} El factor predicho es, por sí solo, la respuesta a la pregunta del apartado~\ref{sec:pregunta}: dice qué fracción de la señal sobrevivirá antes de gastar un ciclo de procesador. La reconstrucción es una aplicación posterior y opcional.
\end{itemize}

\begin{figure}[htbp]
    \centering
    \caption{Pipeline de la propuesta de este trabajo}
    \label{fig:reconstruccion}

    % =====================================================================================
% Donde entra cada dato en la reformulacion — figura propia en TikZ.
%
% ES LA CONTRAFIGURA DE LA 1.2 (cap01/mlqem_flujo). Alli el trazo rojo discontinuo es el
% valor ruidoso medido entrando COMO ENTRADA DEL MODELO, que es lo que obliga a ejecutar
% el circuito antes de poder corregirlo. Aqui el mismo trazo rojo discontinuo entra
% SOLO en la aritmetica final, a la derecha de la linea de puntos que separa el antes del
% despues. Las dos figuras se leen juntas: es el argumento de que la reformulacion no
% rompe el paradigma pre-ejecucion.
%
% CODIGO DE COLOR, el mismo que en el resto de la memoria:
%   azulunir     -> el dato que el modelo SI ve antes de ejecutar
%   acentofig    -> lo que se le HACE (el modelo, la aritmetica)
%   rojoclasico  -> el valor medido, que solo existe DESPUES de ejecutar
%                   (equivale al #b3261e del SVG de la figura 1.2)
%
% ⚠️ La flecha de f-sombrero CRUZA la linea de puntos a proposito: la prediccion nace
%    antes de ejecutar y se usa despues. Lo que NO cruza —y es el punto de la figura— es
%    el trazo rojo, que se queda entero a la derecha.
%
% ⚠️ \shorthandoff{<>} es obligatorio: `babel` en español hace activos < y >, y eso rompe
%    la sintaxis de puntas de flecha de TikZ. Va tambien en el preambulo, pero se repite
%    aqui para que el fichero sea autocontenido si se compila suelto.
% =====================================================================================
\begingroup
\shorthandoff{<>}%
\begin{tikzpicture}[
    >={Stealth[length=2mm]},
    line width=0.6pt,
    font=\footnotesize,
    caja/.style   = {draw=azulunir, line width=0.8pt, rounded corners=2pt,
                     align=center, inner sep=4pt, minimum height=8mm, minimum width=21mm},
    opera/.style  = {draw=acentofig, fill=acentofig!8, line width=0.9pt,
                     rounded corners=2pt, align=center, inner sep=5pt,
                     minimum height=9mm},
    medido/.style = {draw=rojoclasico, line width=0.8pt, rounded corners=2pt,
                     align=center, inner sep=4pt, minimum height=8mm},
  ]

  % --- Columna 1: lo que el modelo ve, y es TODO lo que ve --------------------------
  \node[caja] (circ) at (0,0.85)  {Circuito\\transpilado};
  \node[caja] (tele) at (0,-0.85) {Telemetría\\del chip};

  % --- Columna 2: el modelo ---------------------------------------------------------
  \node[opera, minimum width=15mm] (modelo) at (3.3,0) {Modelo};

  \draw[->, azulunir] (circ.east) -- (modelo.west);
  \draw[->, azulunir] (tele.east) -- (modelo.west);

  % --- Columna 3: la aritmetica de la correccion ------------------------------------
  \node[opera, minimum width=12mm] (arit) at (8.0,0) {$\displaystyle \frac{r}{\hat f}$};

  % La prediccion cruza la frontera: nace antes de ejecutar y se consume despues.
  \draw[->, acentofig] (modelo.east) -- node[above, midway, text=black] {$\hat f$}
        (arit.west);

  \draw[->] (arit.east) -- ++(0.8,0)
        node[right, align=left, inner sep=2pt] {$\langle O\rangle$\\mitigado};

  % --- El valor medido: entra SOLO aqui, y desde el lado de «despues» ----------------
  \node[medido] (med) at (8.0,-2.0) {$\langle O\rangle_{\text{ruidoso}} = r$\\medido al ejecutar};
  \draw[->, rojoclasico, dashed, line width=0.9pt] (med.north) -- (arit.south);

  % --- La frontera: antes / despues de ejecutar --------------------------------------
  \draw[dotted, gray, line width=0.9pt] (6.1,-2.85) -- (6.1,1.75);
  \node[gray, above] at (3.0,1.6)  {antes de ejecutar};
  \node[gray, above] at (8.4,1.6)  {después de ejecutar};

\end{tikzpicture}
\endgroup


    {\footnotesize Fuente: elaboración propia.}
\end{figure}


% =====================================================================================
\section{El umbral de señal}
\label{sec:umbral}
% =====================================================================================
% 🔴 REGLA ANTI-FUGA. Todo este apartado esta medido SOLO sobre `train`. Una critica
%    externa (ago-2026) señalo, con razon, que elegir el umbral mirando validacion o
%    prueba es fuga: el umbral es una constante estimada a partir de datos como cualquier
%    otra. El barrido se rehizo. No reponer aqui ninguna cifra de val ni de test salvo la
%    comprobacion de huerfanas del final, que es un control y no interviene en la eleccion.

El factor de supervivencia es un cociente, y cuando el denominador se acerca a cero se dispara.

\begin{equation}
\label{eq:divergencia_factor}
f = \frac{\langle O \rangle_{\text{ruidoso}}}{\langle O \rangle_{\text{exacto}}} \implies \lim_{|\langle O \rangle_{\text{exacto}}| \to 0} |f| = \infty
\end{equation}

Por lo tanto, surge la necesidad de acotar el denominador.

Para la elección de la cota del valor exacto se ha realizado un barrido previo de un secuencia de umbrales hasta decidir en función de la pérdida de datos frente a la introducción de sesgo por valores muy altos de $f$.

Como se observan en la Tabla~\ref{tab:barrido} y la Figura~\ref{fig:umbral_barrido}, cuanto más alto es el umbral, menos etiquetas quedan y más circuitos se quedan huérfanos (circuitos sin señal en ninguno de los cinco observables), pero más limpia queda la cola del objetivo.

\begin{table}[htbp]
    \centering
    \caption{El barrido del umbral de señal, medido sobre el conjunto de entrenamiento.}
    \label{tab:barrido}
    {\footnotesize
    \begin{tabular}{@{}lrrrrrr@{}}
        \toprule
        \textbf{Umbral} & \textbf{Etiquetas} & \textbf{Huérfanas}
            & \textbf{Huérf. $n{=}11$} & \textbf{$f<0$} & \textbf{$f>1$}
            & \textbf{máx. $|f|$} \\
        \midrule
        $0{,}0005$ & $64{.}737$ & $0{,}00$\,\% & $0{,}00$\,\% & $2{,}50$\,\% & $8{,}64$\,\% & $16{,}45$ \\
        $0{,}0010$ & $63{.}832$ & $0{,}00$\,\% & $0{,}00$\,\% & $2{,}09$\,\% & $8{,}31$\,\% & $11{,}87$ \\
        $\bm{0{,}0020}$ & $\bm{62{.}125}$ & $\bm{0{,}00}$\textbf{\,\%} & $\bm{0{,}00}$\textbf{\,\%}
                        & $\bm{1{,}41}$\textbf{\,\%} & $\bm{7{,}85}$\textbf{\,\%} & $\bm{6{,}67}$ \\
        $0{,}0050$ & $57{.}647$ & $0{,}08$\,\% & $0{,}47$\,\% & $0{,}58$\,\% & $6{,}86$\,\% & $4{,}31$ \\
        $0{,}0100$ & $51{.}006$ & $0{,}93$\,\% & $4{,}43$\,\% & $0{,}16$\,\% & $6{,}03$\,\% & $2{,}29$ \\
        $0{,}0500$ & $23{.}943$ & $24{,}33$\,\% & $48{,}13$\,\% & $0{,}00$\,\% & $4{,}06$\,\% & $1{,}77$ \\
        \bottomrule
    \end{tabular}}

    {\footnotesize Fuente: elaboración propia. Sin umbral hay $66{.}720$ etiquetas. En negrita,
    el umbral elegido.}
\end{table}

% ⚠️ FORMATO: titulo arriba y fuente abajo (Instrucciones §1.3, p. 6).
% Figura propia. La genera umbral.dibujar_barrido_memoria() en el cuaderno
% notebooks/01_eda/01d, y se copia desde
% TFM-Quantum/figures/eda/umbral/04_barrido_memoria.pdf.
%
% El umbral elegido va MARCADO con una vertical en LOS DOS paneles, porque la decision se
% toma mirandolos a la vez: a su izquierda no se gana limpieza, y a su derecha se empieza a
% pagar en muestras huerfanas.
% ⚠️ Va SIN titulo dentro de la imagen: lo pone el \caption, que es lo que pide la norma.
%    Los rotulos de cada panel SI se quedan, porque no titulan la figura.
% ⚠️ Su `figsize` esta calculado para la caja de 15 cm: a \textwidth se imprime a escala
%    ~0,6, y por eso sus cuerpos de letra son de 11-12 pt en origen. Agrandar la figura sin
%    subir las fuentes deja la letra ilegible en el PDF.
\begin{figure}[htbp]
    \centering
    \caption{Lo que cuesta y lo que limpia cada umbral. Eje de abscisas en escala logarítmico.}
    \label{fig:umbral_barrido}

    \includegraphics[width=\textwidth]{cap05/umbral_barrido}

    {\footnotesize Fuente: elaboración propia.}
\end{figure}

El \textbf{umbral} fijado es $\bm{0{,}002}$, con sus consecuencias:

\begin{itemize}
    \item \textbf{Conserva el $93{,}1$\,\% de las etiquetas} ($62{.}125$ de $66{.}720$).
    \item \textbf{Ningún circuito se queda huérfano} en entrenamiento ni en validación, y en el conjunto de
          prueba solo el $0{,}2$\,\%.
    \item \textbf{Limpia la mayor parte de los valores máximos altos}: los factores
          negativos bajan al $1{,}41$\,\% y la cola se recorta de un máximo de $16{,}45$ a uno
          de $6{,}67$. Un factor negativo significa que el valor medido (el ruidoso) cambió de signo respecto
          al exacto. Su justificación aparece en el siguiente anexo~\ref{anexo:factor_negativo}.
    \item \textbf{No penaliza a los circuitos grandes}: con este umbral el porcentaje de huérfanos es nulo en todos los tamaños.
    \item \textbf{Afecta de forma diferente a cada observable}
\end{itemize}

% ⚠️ FORMATO: titulo arriba y fuente abajo (Instrucciones §1.3, p. 6).
% Figura propia. La genera nulos.dibujar_perfil_memoria() en el cuaderno
% notebooks/01_eda/01f, y se copia desde figures/eda/nulos/02_huecos_memoria.pdf.
% ⚠️ Va SIN titulo dentro de la imagen: lo pone el \caption. El umbral, que en la version
%    del cuaderno iba dentro del titulo, baja al pie de figura.
% ⚠️ Su `figsize` esta calculado para la caja de 15 cm: a \textwidth se imprime a escala
%    0,76. Agrandarla sin subir las fuentes deja la letra ilegible en el PDF.
\begin{figure}[htbp]
    \centering
    \caption{Qué porcentaje de sus casillas pierde cada observable con el umbral aplicado.}
    \label{fig:umbral_huecos}

    \includegraphics[width=\textwidth]{cap05/umbral_huecos}

    {\footnotesize Fuente: elaboración propia. Conjunto de entrenamiento, umbral $0{,}002$.}
\end{figure}

Aclaración importante: 

\begin{center} 
      \doublebox{
      \begin{minipage}{0.85\linewidth} \vspace{0.5em}
             \centering 
             \textit{El \textbf{umbral = }$\bm{0{,}002}$ se aplica tras predecir a la hora de evaluar, y no antes, reportando el número de circuitos descartados por la cota.} \vspace{0.5em} \end{minipage}} 
\end{center}


% =====================================================================================
\section{Anatomía de la variable objetivo}
\label{sec:anatomia}
% =====================================================================================

\subsection{Rango}
\label{sec:rango}

\begin{table}[htbp]
    \centering
    \caption{El factor de supervivencia por observable, con el umbral de $0{,}002$ aplicado.}
    \label{tab:anatomia}
    {\footnotesize
    \begin{tabular}{@{}lrrrrrrr@{}}
        \toprule
        \textbf{Observable} & \textbf{Etiquetas} & \textbf{Mín.} & \textbf{Q1}
            & \textbf{Mediana} & \textbf{Q3} & \textbf{Máx.} & \textbf{$f>1$} \\
        \midrule
        Media de $Z$        & $12{.}780$ & $-5{,}19$ & $0{,}65$ & $0{,}87$ & $0{,}97$ & $5{,}21$ & $11{,}05$\,\% \\
        Media de $X$        & $12{.}640$ & $-3{,}05$ & $0{,}61$ & $0{,}83$ & $0{,}96$ & $5{,}21$ & $9{,}87$\,\% \\
        Media de $Y$        & $12{.}672$ & $-5{,}88$ & $0{,}60$ & $0{,}83$ & $0{,}96$ & $4{,}64$ & $9{,}86$\,\% \\
        Paridad             & $11{.}413$ & $-3{,}16$ & $0{,}44$ & $0{,}65$ & $0{,}81$ & $5{,}03$ & $1{,}31$\,\% \\
        Correlador vecinos  & $12{.}620$ & $-3{,}54$ & $0{,}56$ & $0{,}79$ & $0{,}94$ & $6{,}67$ & $6{,}47$\,\% \\
        \bottomrule
    \end{tabular}}

    {\footnotesize Fuente: elaboración propia. Conjunto de entrenamiento.}
\end{table}

% ⚠️ FORMATO: titulo arriba y fuente abajo (Instrucciones §1.3, p. 6).
% Figura propia. La genera factor.dibujar_distribuciones_memoria() en el cuaderno
% notebooks/01_eda/01e, y se copia desde
% TFM-Quantum/figures/eda/factor/06_distribuciones_memoria.pdf.
%
% El eje esta recortado a [-0,5 · 1,5] para que se vea la forma; ese recorte cubre el
% 98,8 % de las etiquetas, y el aviso va DENTRO de la figura, en el hueco que deja la
% rejilla de 2x3. La cola completa esta en la tabla de al lado.
% ⚠️ Va SIN titulo dentro de la imagen: lo pone el \caption, que es lo que pide la norma.
%    Los rotulos de cada panel SI se quedan, porque no titulan la figura.
% ⚠️ Su `figsize` esta calculado para la caja de 15 cm: a \textwidth se imprime a escala
%    ~0,6, y por eso sus cuerpos de letra son de 11-12 pt en origen. Agrandar la figura sin
%    subir las fuentes deja la letra ilegible en el PDF.
\begin{figure}[htbp]
    \centering
    \caption{La distribución del factor de supervivencia en los cinco observables.}
    \label{fig:f_distribuciones}

    \includegraphics[width=\textwidth]{cap05/f_distribuciones}

    {\footnotesize Fuente: elaboración propia.}
\end{figure}

Las cinco distribuciones tienen la misma forma (Figura~\ref{fig:f_distribuciones}): la moda
en uno, una cola larga hacia la izquierda que alcanza valores negativos y una cola corta por encima de uno.

La paridad es la excepción, y en las cinco columnas de la tabla. Su mediana es $0{,}65$
frente al $0{,}83$--$0{,}87$ del resto de observable y solo el
$1{,}31$\,\% de sus etiquetas supera el uno, frente al $11$\,\% de la media de $Z$. Es decir, es el observable que más sufre y el que menos se amplifica.

\subsection{Las dos cotas que impone la física}
\label{sec:cotas}
% ⚠️ APARTADO BREVE A PROPOSITO: es una derivacion, no una discusion. Manda la formula.
%    La regla que sale de la cota inferior y su aplicacion al predecir NO van aqui: son
%    protocolo de evaluacion y viven en §\ref{sec:cota}.
% Fuente: notebooks/01_eda/01e, secciones 2.1 a 2.4.

Las dos salen del mismo sitio: el valor esperado de un observable de Pauli vive en $[-1,1]$
(apartado~\ref{sec:cinco_observables}).

\textbf{Cota superior.} Como $|\langle O\rangle_{\text{ruidoso}}| \le 1$, de la definición del factor se sigue
\begin{equation}
\label{eq:cota_superior}
|f| = \frac{|\langle O\rangle_{\text{ruidoso}}|}{|\langle O\rangle_{\text{exacto}}|} \le \frac{1}{|\langle O\rangle_{\text{exacto}}|} .
\end{equation}

\textbf{Cota inferior.} Sale de desarrollar la reconstrucción de la
ecuación~\eqref{eq:reconstruccion}:

\begin{equation}
\label{eq:mitigado_desarrollo}
\begin{split}
\langle O\rangle_{\text{mitigado}} &= \langle O\rangle_{\text{ruidoso}} + \Delta = \langle O\rangle_{\text{ruidoso}} + \langle O\rangle_{\text{ruidoso}} \left(\frac{1-f}{f}\right) \\
&= \langle O\rangle_{\text{ruidoso}} \left(1 + \frac{1-f}{f}\right) = \frac{\langle O\rangle_{\text{ruidoso}}}{f} .
\end{split}
\end{equation}

\noindent Es decir, \textbf{el valor mitigado es el medido dividido por el factor}. Y como ese
valor mitigado es también un valor esperado de Pauli, tiene que caber en el mismo intervalo:
\begin{equation}
\label{eq:cota_inferior}
-1 \le \langle O\rangle_{\text{mitigado}} \le 1
\quad \Longleftrightarrow \quad \left|\frac{\langle O\rangle_{\text{ruidoso}}}{f}\right| \le 1
\quad \Longleftrightarrow \quad |f| \ge |\langle O\rangle_{\text{ruidoso}}| .
\end{equation}

\noindent Queda así localizada la cota inferior, y es \textbf{una por medición}: no es una
constante, sino un suelo distinto para cada circuito y cada observable.

\noindent \textbf{Ninguna de las $62{.}125$ etiquetas viola ninguna de las dos}, y se acerca a $|f| = |\langle O\rangle_{\text{ruidoso}}|$ solo en los circuitos con valor exacto aproximadamente uno. 

La cota inferior debe aplicarse como un filtro sobre la \emph{predicción}, pues el modelo no garantiza resultados físicos: con $\langle O\rangle_{\text{ruidoso}} = 0{,}01$, predecir $0{,}05$ genera un valor razonable ($0{,}20$), pero $0{,}001$ da $10$, un absurdo físico (ya que deberia de estar contenido en el intervalo $[-1,1]$).

\subsection{El factor de supervivencia frente al tamaño}
\label{sec:f_tamano}

El conjunto de prueba se plantea para extrapolar de once a quince qubits, así que interesa cómo se
mueve la etiqueta con el tamaño dentro del rango que sí se entrena
(Figura~\ref{fig:f_vs_tamano}). Se concluye:

% Figura propia, generada por notebooks/01_eda/01e y copiada desde
% TFM-Quantum/figures/eda/factor/02_f_vs_tamano.pdf. Solo llega a n=11 porque es el tope
% de `train`: mas alla no hay entrenamiento y mirarlo seria mirar el conjunto de prueba.
\begin{figure}[htbp]
    \centering
    \caption{Mediana del factor de supervivencia frente al número de qubits, en el conjunto de
             entrenamiento.}
    \label{fig:f_vs_tamano}

    \includegraphics[width=0.78\textwidth]{cap05/f_vs_tamano}

    {\footnotesize Fuente: elaboración propia.}
\end{figure}

\begin{itemize}
    \item Los cuatro observables promediados sobreviven algo mejor cuanto más grande es el circuito y se sufre mas del fenómeno de amplificación.
    \item Sin embargo, la paridad se derrumba con el aumento de qubits.
\end{itemize}

La consecuencia de esto es que los cinco objetivos son muy distintos. Entonces, al describirlos y reportar resultados es necesario tratarlos de forma independiente.

\subsection{Los cinco objetivos entre sí}
\label{sec:correlacion_objetivos}

Se planteó la siguiente pregunta: ¿un modelo por observable, o uno solo con cinco salidas? Un tronco compartido solo tiene sentido si los cinco tienen algún tipo de correlación.

% ⚠️ FORMATO: titulo arriba y fuente abajo (Instrucciones §1.3, p. 6).
% Figura propia. La genera factor.dibujar_correlaciones_memoria() en el cuaderno
% notebooks/01_eda/01e, y se copia desde figures/eda/factor/05_correlaciones_memoria.pdf.
% ⚠️ Va SIN titulo dentro de la imagen: lo pone el \caption. Los rotulos «Pearson» y
%    «Spearman» SI se quedan, porque no titulan la figura sino cada panel.
% ⚠️ UN MAPA DE COLOR POR PANEL, por decision expresa del autor (sep-2026). El precio es
%    que los dos paneles ya no se comparan por el color; los dos conservan el dominio
%    [-1, 1] para que la comparacion numerica siga siendo directa.
% ⚠️ La escala llega a -1 aunque ningun valor baje de 0,40. Es deliberado: es lo que enseña
%    que las diez correlaciones son positivas y que ninguna roza el cero.
\begin{figure}[htbp]
    \centering
    \caption{Correlación con signo entre los cinco objetivos, medida de dos formas. Cada
             panel lleva su propia escala de color; las dos van de $-1$ a $+1$.}
    \label{fig:f_correlaciones}

    \includegraphics[width=\textwidth]{cap05/f_correlaciones}

    {\footnotesize Fuente: elaboración propia.}
\end{figure}

La correlación de Spearman oscila entre $0{,}63$ y $0{,}75$, y la de Pearson entre $0{,}4$ y $0{,}6$. Estos datos justifican la elección de un \textbf{tronco común con cinco salidas}.


% =====================================================================================
        % ⭐ resultado propio
\include{capitulos/06_planteamiento_comparativa}

% --- Desarrollo de la comparativa ---------------------------------------------------
\chapter{Desarrollo de la comparativa}
\label{cap:comparativa}

% =====================================================================================
% GUION DEL CAPITULO — ESTRUCTURA FIJADA POR EL AUTOR (sep-2026).
%
% 🔴 NO AÑADIR SECCIONES SIN CONSULTARLO. La estructura acordada es exactamente esta:
%      7.1 Estrategia de entrenamiento   7.3 Random Forest
%      7.2 Ridge                         7.4 GNN
%
% ⚠️ La version anterior del capitulo se DESCARTO entera (sep-2026): duplicaba el
%    capitulo~\ref{cap:planteamiento} —tenia sus propias secciones de particion,
%    preprocesado y metricas— y ademas repetia la explicacion de los tres modelos. Si
%    hace falta recuperar algo, esta en el historial del repositorio.
%
% 🔴 REPARTO CON EL CAPITULO 6, que es la regla que evita volver a duplicar:
%      cap. 6 -> los DATOS y el diseño experimental: de donde salen, que dice la EDA,
%                como se preprocesan y con que metricas se juzgan.
%      cap. 7 -> los MODELOS: cual es cada uno, por que esta, con que hiperparametros y
%                como se ha entrenado.
%    Si una frase de aqui habla del conjunto de datos, del umbral o de una metrica, esta
%    en el capitulo equivocado: se cita el apartado del 6 y se sigue.
%
% ⚠️ EXTENSION: el capitulo NO debe ser largo. Una figura por modelo, y la tabla de las
%    tres convoluciones al final de §7.4.
%
% FUENTES, por seccion:
%   7.1 -> notebooks/03_ridge_rf/NOTAS_ENTRENAMIENTO.md
%          notebooks/04_gnn/NOTAS_ENTRENAMIENTO.md
%          doc/changes/2026-08_cambio_decisiones.md §4 (la escalera de seis peldaños)
%   7.2 -> src/baselines.py · scripts/train_ridge.py
%   7.3 -> src/baselines.py · scripts/train_rf.py · liao2024machine (el mejor tabular)
%   7.4 -> src/gnn_model.py · src/train.py · configs/gnn_config.yaml
%          notebooks/04_gnn/NOTAS_ENTRENAMIENTO.md · doc/changes/2026-08_cambio_decisiones.md §9
%
% 🔴 CIFRAS PROHIBIDAS (son del dataset v1, obsoleto): «30-70 saltos» de distancia entre
%    nodos —la del v2 es 44— y «42,7x» / «45,4x» del QFT —la citable es 3,69x—.
%
% ⚠️ NINGUN RESULTADO ENTRA AQUI. Este capitulo dice como se entrena; que salio de ello
%    es el capitulo~\ref{cap:resultados}. Si aparece un R² o un MAE, esta mal colocado.
% =====================================================================================


% =====================================================================================
\section{Estrategia de entrenamiento}
\label{sec:entrenamiento}
% =====================================================================================

La estrategia del entrenamiento se compone de estos seis peldaños:

\begin{enumerate}
    \item \textbf{No mitigar.} El suelo: si no se bate, el método empeora las cosas y no sirve.
    \item \textbf{Factor de supervivencia constante}, la media de entrenamiento. ¿Aporta saber la media?
    \item \textbf{Ridge con una sola variable}, la duración total. ¿Aporta saber cuánto dura
          el circuito?
    \item \textbf{Ridge con todas las variables.} ¿Aporta solo la linealidad?
    \item \textbf{Random Forest.} ¿Aporta la no linealidad?
    \item \textbf{Red de grafos.} ¿Aporta la estructura del circuito?
\end{enumerate}

Los tres modelos que se comparan son los tres últimos peldaños y cada uno añade una capacidad sobre el anterior: Ridge es el
\textbf{baseline} y fija el suelo, Random Forest añade no \textbf{linealidad} y está porque es el mejor modelo tabular de los que compara la referencia \cite{liao2024machine}. La red de grafos cambia la representación de los datos, y la diferencia entre ella y los dos primeros es \textbf{el valor de la
estructura que la agregación descarta}.

Cabe recalcar una asimetría presente. Ridge y la red de
grafos se entrenan con una \textbf{pérdida de Huber}, que penaliza los errores grandes de forma lineal y evita que los factores de supervivencia extremos se sesgen el entrenamietno. Sin embargo, \textbf{Random Forest no la admite} y no se emplea.


% =====================================================================================
\section{Ridge}
\label{sec:sol_ridge}
% =====================================================================================

Una regresión lineal con penalización $L_2$ \cite{hoerl1970ridge}. No se usan mínimos cuadrados
básico por un motivo medido en el apartado~\ref{sec:colinealidad}: aun después de descartar las
veinte columnas redundantes y se prefiere que el modelo recaba información de ella ajustándose solo.

Se ajusta \textbf{un modelo por observable}, cinco en total, sobre las $34$ columnas
preprocesadas y descartando las filas sin etiqueta de es observable. El único
hiperparámetro es la constante de regularización, barrida en diez valores de $10^{-3}$ a
$10^{6}$. El óptimo cae en $100$ para la media de $Z$ y el correlador de vecinos, y en $1{.}000$
para los otros tres.

La Figura~\ref{fig:ridge_coef} recoge los doce coeficientes de mayor magnitud. Se muestran
\textbf{con su signo}, que es lo único que este modelo da y los otros dos no: la dirección en
la que cada variable empuja al factor de supervivencia.

% ⚠️ FORMATO: titulo arriba y fuente abajo (Instrucciones §1.3, p. 6).
% Figura propia. La genera modelos.dibujar_ridge_coeficientes() y se copia desde
% TFM-Quantum/figures/modelos/ridge_coeficientes.pdf. Sin titulo interno.
\begin{figure}[htbp]
    \centering
    \caption{Los doce coeficientes mayores de Ridge, con su signo, promediando los cinco
             observables.}
    \label{fig:ridge_coef}

    \includegraphics[width=0.92\textwidth]{cap07/ridge_coeficientes}

    {\footnotesize Fuente: elaboración propia.}
\end{figure}

% =====================================================================================
\section{Random Forest}
\label{sec:sol_rf}
% =====================================================================================

Un bosque de árboles de decisión \cite{breiman2001random} añade sobre Ridge no linealidad.

Cabe recalcar que este modelo \textbf{no extrapola} y no será util para comparar el \textbf{eje tamaño}.

Se entrena igual que Ridge (un modelo por observable, las mismas $34$ columnas y los mismos
descartes). El escalado no le afecta, pero \textbf{se le aplican las mismas transformaciones}: no por necesidad, sino para que la única diferencia entre los dos
modelos tabulares sea el modelo. La rejilla fija $300$ árboles y barre la profundidad máxima
($\{$sin límite, $6$, $12$, $24\}$), las hojas mínimas ($\{1, 5, 20, 50\}$) y las variables por
división ($\{\sqrt{\cdot}, 0{,}5, 1{,}0\}$). Los óptimos salen \textbf{poco profundos}: profundidad $6$ en tres
de los cinco observables, con hojas de $5$ o $50$ muestras y la mitad de las variables por
división. En la Figura~\ref{fig:rf_imp} se muestra que el modelo se apoya en muy pocas columnas.

% Figura propia. La genera modelos.dibujar_rf_importancias() y se copia desde
% TFM-Quantum/figures/modelos/rf_importancias.pdf. Sin titulo interno.
\begin{figure}[htbp]
    \centering
    \caption{Las doce variables más importantes para Random Forest, promediando los cinco
             observables.}
    \label{fig:rf_imp}

    \includegraphics[width=0.92\textwidth]{cap07/rf_importancias}

    {\footnotesize Fuente: elaboración propia.}
\end{figure}

Las dos figuras cuentan lo mismo y coinciden con que \textbf{la duración total y la profundidad tiene mayor peso en el modelo},
con un $0{,}39$ y un $0{,}33$ de importancia frente al $0{,}03$ del siguiente.


% =====================================================================================
\section{La red de grafos}
\label{sec:gnn}
% =====================================================================================

Una red de paso de mensajes sobre el grafo del circuito. Es la que más justificación necesita,
porque el diseño se apartó de la arquitectura de referencia del área.

\textbf{Paso de mensajes y no atención.} El trabajo más cercano \cite{bao2025gtraqem} usa un
\textit{graph transformer} \emph{sin} paso de mensajes; sin embargo, el computo disponible era insuficiente y por eso ser recurrió a modelos de grafos menos pesados.

\textbf{El resumen global es la pieza que hace funcionar el modelo}
(Figura~\ref{fig:gnn_arq}). Se escogió incorporar un vector que conecte a todos lo nodos para ahorra la distancia de paso entre cualquier nodo y, así, usar ese vector como atajo. La idea se toma del mismo trabajo\cite{bao2025gtraqem}, con un cambio: allí el
vector es \textbf{unidireccional}, y eso basta en un \textit{transformer} porque su atención ya
conecta todo con todo; sin embargo, en una red de paso de mensajes no, y lo que queda es un \textit{pooling} aprendido calculando con
\textbf{media}. Los grafos entran \textbf{completos}.

% FIGURA PROPIA en TikZ. Codigo en figuras/cap07/gnn_arquitectura_tikz.tex.
% ⚠️ Lo que la figura tiene que dejar claro son las DOS flechas del resumen global: lee y
%    devuelve. Es lo que lo distingue de un pooling aprendido, y en prosa no se ve.
\begin{figure}[!ht]
    \centering
    \caption{La arquitectura. El resumen global lee de todos los nodos y les devuelve: son
             las dos flechas, no una.}
    \label{fig:gnn_arq}

    % =====================================================================================
% La arquitectura del MPNN — figura propia en TikZ.
%
% Lo que tiene que quedar claro y no se entiende sin dibujo: el resumen global es
% BIDIRECCIONAL. Lee de todos los nodos y DEVUELVE a todos los nodos, y esa flecha de
% vuelta es la que lo distingue de un pooling aprendido. Por eso hay dos flechas entre la
% capa de mensajes y el resumen, y no una.
%
% CODIGO DE COLOR, el de la memoria:
%   azulunir  -> el dato (el grafo de entrada y la salida)
%   acentofig -> lo que se le hace (las capas y el resumen)
%
% ⚠️ \shorthandoff{<>} es obligatorio: `babel` en español hace activos < y >.
% =====================================================================================
\begingroup
\shorthandoff{<>}%
\hyphenpenalty=10000\exhyphenpenalty=10000\relax

\begin{tikzpicture}[
    >={Stealth[length=2mm]},
    font=\footnotesize,
    dato/.style  = {draw=azulunir, line width=0.9pt, rounded corners=3pt, fill=azulunir!6,
                    align=center, inner sep=5pt, text width=2.1cm},
    capa/.style  = {draw=acentofig, line width=0.9pt, rounded corners=3pt,
                    fill=acentofig!8, align=center, inner sep=5pt, text width=2.3cm},
    resumen/.style = {draw=acentofig, line width=1.2pt, rounded corners=3pt,
                      fill=acentofig!16, align=center, inner sep=5pt, text width=2.6cm},
  ]

  % --- la cadena principal -------------------------------------------------------------
  \node[dato]  (entrada) at (0,0)
    {\textbf{Grafo}\\{\scriptsize $21$ dims por nodo}};
  \node[capa]  (mp)      at (3.4,0)
    {\textbf{Paso de mensajes}\\{\scriptsize $4$ capas · aristas dirigidas}};
  \node[capa]  (lectura) at (7.2,0)
    {\textbf{Lectura}\\{\scriptsize resumen $\oplus$ \textit{pooling}}};
  \node[dato]  (salida)  at (10.4,0)
    {\textbf{$\hat f$}\\{\scriptsize cinco salidas}};

  \draw[->, azulunir,  line width=1pt] (entrada.east) -- (mp.west);
  \draw[->, acentofig, line width=1pt] (mp.east)      -- (lectura.west);
  \draw[->, azulunir,  line width=1pt] (lectura.east) -- (salida.west);

  % --- el resumen global, encima, con SUS DOS flechas ------------------------------------
  \node[resumen] (virtual) at (3.4,2.0)
    {\textbf{Resumen global}\\{\scriptsize un vector por grafo · pesos propios}};

  \draw[->, acentofig, line width=1pt] ($(mp.north)+(-0.42,0)$) -- ($(virtual.south)+(-0.42,0)$)
        node[midway, left=1pt, font=\scriptsize, text=black] {lee};
  \draw[->, acentofig, line width=1pt] ($(virtual.south)+(0.42,0)$) -- ($(mp.north)+(0.42,0)$)
        node[midway, right=1pt, font=\scriptsize, text=black] {devuelve};

  \draw[->, acentofig, line width=0.8pt, dotted]
        (virtual.east) to[out=0, in=110] (lectura.north);

  % --- la nota que justifica la vuelta ---------------------------------------------------
  \node[align=left, font=\scriptsize, text=gray, text width=3.7cm] at (8.9,2.1)
    {con $4$ capas un nodo alcanza\\el $1{,}25$\,\% de los pares:\\el resumen es la única vía\\al largo alcance};

  % --- lo que envuelve cada capa ---------------------------------------------------------
  \node[align=center, font=\scriptsize, text=gray] at (3.4,-1.35)
    {cada capa: normalización + conexión residual};

\end{tikzpicture}%
\endgroup


    {\footnotesize Fuente: elaboración propia.}
\end{figure}

\textbf{Tres convoluciones, en orden de complejidad creciente} (Tabla~\ref{tab:convs}): se sube
de peldaño solo si el anterior se queda corto, de modo que si la atención local mejora se sabe
que aporta, y si no, también.

\begin{table}[!ht]
    \centering
    \caption{Las tres convoluciones que se prueban y la configuración con la que se entrenó
             cada una.}
    \label{tab:convs}
    {\footnotesize
    \begin{tabular}{@{}p{0.20\textwidth}p{0.27\textwidth}rrr@{}}
        \toprule
        \textbf{Convolución} & \textbf{Qué añade} & \textbf{Oculta} & \textbf{Tasa}
            & \textbf{Parámetros} \\
        \midrule
        \texttt{SAGEConv} \cite{hamilton2017sage}
            & Agrega los vecinos por media. Ligera y estable: es la de arranque
            & $256$ & $10^{-3}$ & $1{,}06$~M \\
        \texttt{GINConv} \cite{xu2019gin}
            & Suma estricta más un perceptrón interno. La más expresiva; distingue mejor
              topologías parecidas
            & $256$ & $3\cdot10^{-4}$ & $1{,}06$~M \\
        \texttt{GATv2Conv} \cite{brody2022gatv2}
            & \textbf{Atención local}: aprende a pesar más a los vecinos con peor telemetría
            & $160$ & $10^{-3}$ & $1{,}04$~M \\
        \bottomrule
    \end{tabular}}

    {\footnotesize Fuente: elaboración propia. Las tres con cuatro capas.}
\end{table}

\textbf{Las tres variante se prueban con el mismo presupuesto de
parámetros}

El resto de la rejilla barre la tasa de aprendizaje ($\{3\cdot10^{-4}, 10^{-3},
3\cdot10^{-3}\}$), la dimensión oculta ($\{128, 192, 256\}$) y el número de capas ($\{4, 6\}$),
con lotes de $32$ grafos y regularización de $10^{-5}$. El entrenamiento se hace en una máquina
con tarjeta gráfica en la nube \cite{google2026cloud}, porque en local no hay ninguna.     % ⏳ la EJECUCIÓN: barrido, hiperparámetros, incidencias

% --- Resultados y discusión ---------------------------------------------------------
\include{capitulos/08_resultados}      % ⏳ BLOQUEADO: resultados Y su discusión

% --- Conclusiones -------------------------------------------------------------------
\chapter{Conclusiones y trabajo futuro}
\label{cap:conclusiones}

% =====================================================================================
% ⚠️ EXIGENCIA EXPLICITA DE LA NORMA (Instrucciones seccion 2.7):
%    «cada objetivo del trabajo se enlazara con una conclusion». Por eso §9.1 recorre los
%    cinco especificos UNO A UNO y cierra con el general. NO fundirlos en un parrafo.
%
% ⚠️ EXTENSION FIJADA POR EL AUTOR (sep-2026): DOS paginas. Esta AJUSTADO al limite: si se
%    añade una linea, §9.4 se cae a una tercera pagina. Recortar antes de añadir.
%
% ⚠️ NINGUNA CIFRA NUEVA. Todas las de este capitulo estan medidas y redactadas en el
%    capitulo~\ref{cap:resultados}. Si aparece un numero que no este alli, esta inventado.
% =====================================================================================

\section{Cumplimiento de los objetivos}
\label{sec:cumplimiento}

Los cinco objetivos específicos del apartado~\ref{sec:objetivos_especificos} se cierran así:

\begin{enumerate}
    \item \textbf{Comprensión de los fundamentos de la computación cuántica.} Cumplido. El
          Capítulo~\ref{cap:fundamentos} fija lo imprescindible: qué es un qubit, qué mide un
          observable y por qué un circuito es un grafo dirigido.

    \item \textbf{Construcción del conjunto de datos.} Cumplido. Veinte mil circuitos de cinco
          a quince qubits, etiquetados con su degradación exacta sobre la calibración real de
          un procesador de IBM durante cuarenta y dos días, y particionados en tres ejes de
          generalización separables (Capítulo~\ref{cap:planteamiento}).

    \item \textbf{Planteamiento y análisis de la variable objetivo.} Cumplido, y con el
          hallazgo que reorientó el trabajo: la degradación con signo \textbf{no es predecible
          antes de ejecutar}. La reformulación sobre el
          \textbf{factor de supervivencia} sí es aprendible (apartado~\ref{sec:reformulacion}).

    \item \textbf{Realización de la comparativa.} Cumplido, con un matiz sobre el `modelo
          óptimo' que se pedía conseguir: \textbf{no existe uno estrictamente mejor}. Random
          Forest gana en la métrica que decide y es el único que nunca empeora el observable al
          corregirlo, pero las redes de grafos lo superan por márgenes grandes en dos
          situaciones (Capítulo~\ref{cap:resultados}).

    \item \textbf{Cuantificación del aporte de la estructura.} Cumplido, y la respuesta es en
          buena parte negativa. El grafo solo gana donde hay que salirse de lo visto: la media
          de $X$ al extrapolar a catorce y quince qubits (apartado~\ref{sec:res_tamano}) y algunos poco circuitos en observables concretos(apartado~\ref{sec:res_tipo}). En el régimen conocido,
          treinta y cuatro columnas agregadas bastan.
\end{enumerate}

El \textbf{objetivo general} (apartado~\ref{sec:objetivo_general}) se cumple. La
\textbf{hipótesis nula} se refuta en su segunda mitad (los cinco modelos baten con holgura la
línea base del factor constante) y solo en parte en la primera: hay diferencias medibles entre
modelos, pero ninguno domina los tres ejes a la vez.

\section{Aportaciones}
\label{sec:aportaciones}

\begin{enumerate}
    \item \textbf{Un conjunto de datos que no existía.} Su valor no está en el tamaño, sino en
          los tres ejes, que permiten medir por separado tres formas distintas de fallar fuera
          de distribución.

    \item \textbf{La demostración medida de que el objetivo natural del problema no es
          predecible antes de ejecutar}, con su explicación y su reformulación.

    \item \textbf{Una comparativa que separa dos preguntas que suelen confundirse}: cuánto
          acierta un modelo sobre el factor, y cuánto mejora el observable al aplicarlo. La
          segunda no se sigue de la primera, y aquí se mide en lugar de suponerse
          (apartado~\ref{sec:res_mitigacion}).
\end{enumerate}

\section{Limitaciones}
\label{sec:limitaciones_finales}
% NO minimizarlas: declararlas hace creible el resto.

\begin{itemize}
    \item El modelo de ruido \textbf{no incluye crosstalk}: las conclusiones son válidas dentro
          de ese modelo.
    \item No hay validación en \textbf{hardware real}: el plan abierto de IBM da $10$ minutos
          de QPU al mes \cite{ibm2026productos}.
    \item El \textit{ground truth} exige simulación exacta, lo que \textbf{acota el tamaño} de
          los circuitos a quince qubits.
    \item El transpilador de Qiskit \textbf{no es determinista}, así que la reproducibilidad
          del conjunto es \textbf{del plan de generación}.
    \item El ruido se modela como factor multiplicativo puro, sin término afín
          (Anexo~\ref{anexo:factor_negativo}).
\end{itemize}

\section{Trabajo futuro}
\label{sec:futuro}
% Fuente: TFM-Quantum/IDEAS_FUTURAS.md

\begin{itemize}
    \item \textbf{Validación en hardware real}: la única limitación que no se resuelve con más
          cómputo clásico.
    \item \textbf{Incorporar el crosstalk} al modelo de ruido y volver a medir los tres ejes.
    \item \textbf{Modelar el término afín}, para los circuitos en los que el factor cambia de signo.
    \item \textbf{Ampliar la batería de observables}, hoy limitada a cinco.
\end{itemize}
    % ⏳ BLOQUEADO: depende del capítulo 8

% =====================================================================================
% Bibliografía
% =====================================================================================
% ⚠️ Sin \nocite{*}: solo deben aparecer las obras REALMENTE citadas en el texto.
% Incluir referencias no citadas es un error que un tribunal detecta.
\bibliographystyle{apacite}
\bibliography{bibliografia}

% =====================================================================================
% Los anexos NO computan a efectos de extensión (Instrucciones §2.9).
\appendix

% Va el primero de los anexos por ser una declaración exigida por la normativa.
\chapter{Declaración de uso de herramientas de inteligencia artificial}
\label{anexo:ia}

% =====================================================================================
% ⚠️ ANEXO OBLIGATORIO SI SE HAN USADO HERRAMIENTAS DE IA.
%
% «Guía del Trabajo Fin de Estudios», apartado 4.4 «Actividades y entregas», PAGINA 22:
%
%   «...en caso de usar alguna de estas herramientas, el director de TFE debe validar
%    su uso adecuado y ético según las recomendaciones de la universidad. Igualmente,
%    dicho uso debe quedar explícitamente reflejado en el documento incluyendo las
%    referencias específicas a las herramientas de IA utilizadas según las normas de
%    citación oportunas. En el caso de que la universidad detecte un uso inadecuado de
%    herramientas de IA, o una citación inadecuada de las mismas, podrá NO AUTORIZAR el
%    trabajo y declararlo no apto para su defensa. En este caso, podrás perder la
%    convocatoria en curso.»
%
% 🔴 DOS COSAS PENDIENTES, Y LAS DOS SON BLOQUEANTES:
%    1. VALIDACION DEL DIRECTOR. La norma la exige de forma explicita. Hay que
%       planteárselo en sesión y dejar constancia.
%    2. REFERENCIAS APA de cada herramienta en bibliografia.bib, y citarlas desde aquí.
%       Estan añadidas pero SIN VERIFICAR: falta confirmar el nombre exacto del producto,
%       la version y el periodo de uso. No inventarlos.
%
% ⚠️ La declaracion tiene que ser COMPLETA. Omitir una categoria de uso es justo lo que
%    la norma castiga. Si aparece un uso nuevo, se añade aqui.
% =====================================================================================

\section{Alcance de la declaración}
\label{sec:ia_alcance}
% Parrafo de entrada: por que existe este anexo y a que se compromete.
Este anexo detalla, de forma exhaustiva, en qué tareas se han empleado herramientas de
inteligencia artificial generativa durante la realización del trabajo, en cumplimiento de
lo establecido en la guía de la asignatura.

\section{Herramientas empleadas y uso concreto}
\label{sec:ia_herramientas}

\begin{table}[htbp]
    \centering
    \begin{tabular}{@{}p{0.28\textwidth}p{0.62\textwidth}@{}}
        \toprule
        \textbf{Herramienta} & \textbf{Uso} \\
        \midrule
        Claude Code (Anthropic)
            & Desarrollo del código de generación y simulación de circuitos cuánticos
              con Qiskit, y del resto de código ajeno al aprendizaje automático. \\[0.4em]
        Claude Code (Anthropic)
            & Generación de las figuras e ilustraciones de la memoria. \\[0.4em]
        Claude Code (Anthropic)
            & Asistencia en la redacción, la estructuración y la revisión del documento. \\[0.4em]
        Google Gemini
            & Búsqueda bibliográfica en profundidad para la elaboración del estado del
              arte. \\
        \bottomrule
    \end{tabular}

    \caption{Herramientas de inteligencia artificial empleadas y tareas en las que se han
             utilizado.}
    \label{tab:ia}

\end{table}

% 🔴 LAS DOS HERRAMIENTAS SE CITAN AQUI CON \cite{}, y es OBLIGATORIO: la norma pide «las
%    referencias específicas a las herramientas de IA utilizadas según las normas de
%    citación oportunas». Una tabla sola no basta. Las entradas del .bib siguen el formato
%    APA 7 para IA generativa: autor, año, producto y versión, descriptor entre corchetes
%    y URL.

Las dos herramientas quedan referenciadas según la normativa APA que exige la titulación:
\textbf{Claude Code} \cite{anthropic_claude} y \textbf{Gemini Pro}
\cite{google_gemini}.

Las figuras de la memoria son de \textbf{dos clases distintas}, y su pie lo refleja. Las
gráficas que representan medidas del trabajo se han producido con código propio sobre los
datos generados para esta memoria, y su fuente es elaboración propia; los \textbf{diagramas
conceptuales}, que no representan ningún dato, se han elaborado con Claude Code y llevan esa
herramienta como fuente.

% \chapter{Código fuente y datos analizados}
% ⚠️ Instrucciones §2.9: «Es recomendable que el estudiante incluya en su memoria la URL
%    del repositorio donde tiene alojado el código fuente desarrollado durante el TFE. El
%    estudiante debe ser el único autor del código y único propietario del repositorio.
%    En el repositorio no debe haber commit de ningún otro usuario del repositorio.»
% PENDIENTE: URL del repositorio, una vez esté limpio.

% Sale del capitulo 5 (§5.3 y §5.5 remiten aqui) para no desviar la linea de trabajo.
\chapter{El término afín del ruido}
\label{anexo:factor_negativo}

% =====================================================================================
% ANEXO CORTO A PROPOSITO. Sale del capitulo 5 para no desviar la linea de trabajo: alli
% §5.5 solo declara la simplificacion y remite aqui, y §5.3.2 apunta aqui al hablar de los
% factores negativos.
%
% 🔴 ES EL UNICO SITIO DE LA MEMORIA donde se desarrolla el modelo afin. No reponer la
%    formula ni el descarte de `b` en §5.5: quedaria dicho dos veces.
%
% Fuente de las cifras: doc/changes/2026-08_cambio_decisiones.md §1 y §3 (dataset v2,
% ajuste sobre una submuestra de 3.000 circuitos). Todas citables.
%
% ⚠️ El paso de `f = f_real + b/e` es algebra, no una medicion: sale de dividir la
%    ecuacion afin por `e`. No hace falta respaldarlo con ninguna cifra.
% =====================================================================================

El factor de supervivencia que define el capítulo~\ref{cap:variable} asume que el ruido es
\textbf{multiplicativo puro}: que el valor medido es el exacto multiplicado por una fracción
que sobrevive. Medido, no lo es del todo. El modelo real es \textbf{afín},
\begin{equation}
\label{eq:modelo_afin}
\langle O \rangle_{\text{ruidoso}} = f \cdot \langle O \rangle_{\text{exacto}} + b ,
\end{equation}
\noindent donde $b$ es un desplazamiento que \textbf{no depende de cuánta señal había}. Tiene
explicación física: la relajación hace caer los qubits hacia $\ket{0}$, cuyo valor de $Z$ es
$+1$, así que el ruido empuja la base $Z$ hacia valores positivos aunque el valor exacto sea
cero.

\section{Por qué el factor se dispara, y por qué cambia de signo}

La consecuencia se ve dividiendo la ecuación~\eqref{eq:modelo_afin} por el valor exacto, que
es lo que hace la definición del factor. Escribiendo $e$ para el valor exacto:
\begin{equation}
\label{eq:factor_afin}
\frac{\langle O \rangle_{\text{ruidoso}}}{e} = f + \frac{b}{e} .
\end{equation}

\noindent El factor que se mide no es $f$, sino $f$ \textbf{más un término que crece cuando el
denominador se encoge}. Con un valor exacto grande, $b/e$ es despreciable y las dos
formulaciones coinciden. Con un valor exacto diminuto (que es el caso típico en este conjunto
de datos, donde la mediana del valor absoluto ronda las tres centésimas) ese término domina, y
si además $b$ tiene el signo contrario al de $e$, \textbf{el cociente entero se vuelve
negativo}. De ahí salen los factores negativos del apartado~\ref{sec:umbral}.

\section{Por qué no se modela}

Para simplificar el trabajo.

% \chapter{Catálogo completo de las variables agregadas}
% PENDIENTE.

% \chapter{Tablas de resultados desglosadas}
% PENDIENTE. Aquí van las tablas largas que no caben en el capítulo 10.

% =====================================================================================
% Índice de acrónimos — opcional según Instrucciones §2.10, pero muy recomendable en
% este trabajo: NISQ, QEC, QEM, ZNE, PEC, M3, GEM, QCR, DAG, OOD, GNN, MAE, RMSE, QPU…
% =====================================================================================
% Sin numerar: es un apartado propio (§2.10), no un anexo más.
% \chapter*{Índice de acrónimos}
\addcontentsline{toc}{chapter}{Índice de acrónimos}
% PENDIENTE.

\end{document}
